\documentclass[a4paper,10pt]{article}
\usepackage[utf8]{inputenc}
\usepackage[T2A]{fontenc}
\usepackage[russian]{babel}
\usepackage{amsmath, amssymb, amsfonts}
\usepackage{graphicx}
\usepackage{geometry}
\usepackage{xcolor}
\usepackage{hyperref}
\usepackage{cite}
\usepackage{float}
\usepackage{subcaption} % если нужны подрисуночные подписи

\geometry{left=1.5cm, right=1.5cm, top=1.5cm, bottom=1.5cm}

% ============================================================
% ЦВЕТОВОЕ ВЫДЕЛЕНИЕ МОДИФИКАЦИЙ
% Измените цвет здесь для всех модификаций сразу
% ============================================================
\definecolor{modcolor}{rgb}{0.0, 0.55, 0.0}  % зелёный (измените по желанию)
% \definecolor{modcolor}{rgb}{0.0, 0.0, 0.7}  % синий (альтернатива)
\definecolor{redcolor}{rgb}{0.7, 0.0, 0.0}  % красный (альтернатива)

\newcommand{\redtext}[1]{\textcolor{redcolor}{#1}}
\newcommand{\redmath}[1]{\textcolor{redcolor}{#1}}
\newcommand{\redblock}[1]{\colorbox{redcolor!10}{\parbox{\dimexpr\linewidth-2\fboxsep}{#1}}}

\newcommand{\modtext}[1]{\textcolor{modcolor}{#1}}
\newcommand{\modmath}[1]{\textcolor{modcolor}{#1}}
\newcommand{\modblock}[1]{\colorbox{modcolor!10}{\parbox{\dimexpr\linewidth-2\fboxsep}{#1}}}


\title{\textbf{}}
\author{}
\date{\today}

\begin{document}



% ============================================================
% §17. ПОНДЕРОМОТОРНЫЕ СИЛЫ (стр. 88–91)
% ============================================================

\section*{\S\,17. Пондеромоторные силы.}

% стр. 88 (левая колонка — продолжение §16)

ников расположены на их поверхностях, так что $\rho = 0$, и что
потенциал каждого проводника постоянен на всём его протяжении,
получим из (15.6):
$$
    W = \frac{1}{2}\sum_{i=1}^n\oint_{S_i}\sigma\varphi\,dS
      = \sum_{i=1}^n\varphi_i\oint_{S_i}\sigma\,dS.
$$
Интеграл $\sigma$ по поверхности проводника равен его общему заряду
$e_i$, поэтому
\begin{equation}
    W = \frac{1}{2}\sum_{i=1}^n e_i\varphi_i.
    \tag{16.7}
\end{equation}

Эту формулу, выражающую \textit{полную} энергию \textit{заряжённых
проводников}, не нужно смешивать с вполне аналогичной формулой (15.4),
выражающей \textit{взаимную} энергию \textit{точечных зарядов}; в
последней формуле $\varphi_i$ не является полным потенциалом поля
в местонахождении заряда $e_i$ (см.\ §\,15). Заметим, что выражение
(15.9) энергии конденсатора является частным случаем формулы (16.7)\,\footnote{Предлагаем
читателю в качестве упражнения непосредственным вычислением убедиться
на конкретных примерах плоского, цилиндрического и шарового конденсаторов
в тождестве значений энергии, получаемых по формуле (15.9), с одной
стороны, и по формуле (16.2) — с другой.}.

\textbf{Задача 14.} Показать, что (собственная) электрическая энергия
заряженного шара радиуса $a$ равна
$$
    W = \frac{e^2}{2a},
$$
если заряд $e$ распределён по поверхности шара (проводник), и равна
$$
    W = \frac{3e^2}{5a},
$$
если заряд равномерно распределён по всему объёму шара.

\textbf{Задача 15.} Показать, что если расстояние между зарядами $e_1$
и $e_2$ достаточно велико по сравнению с их размерами (точечные заряды),
то выражение (16.6) взаимной энергии $W_{12}$ этих зарядов сводится
к выражениям (15.2) и (15.3).

\bigskip
\noindent\textbf{\S\,17. Пондеромоторные силы.}

\medskip
\textbf{1.} По данному в §\,2 определению, основная величина,
характеризующая электрическое поле, — его напряжённость $\mathbf{E}$ —
равна рассчитанной на единицу заряда пондеромоторной (механической)
силе, действующей в данной точке поля на пробный заряд (заряженное

% стр. 89 (правая колонка)

тело). При этом нужно, конечно, иметь в виду, что сила, действующая
на какой-либо заряд, определяется, очевидно, напряжённостью того поля,
в которое помещён этот заряд, а не того поля, которое возбуждается им
самим. Следовательно, говоря, например, что сила, действующая на
точечный заряд $e$, равна:
\begin{equation}
    \mathbf{F} = e\mathbf{E}
    \tag{17.1}
\end{equation}
[ср.\ уравнение (2.2)], мы должны понимать под $\mathbf{E}$ напряжённость
поля, возбуждаемого всеми зарядами системы, \textit{кроме} самого
заряда $e$.

\textbf{2.} Вопрос о силе, действующей на \textit{поверхностные заряды},
требует специального исследования, ибо напряжённость поля $\mathbf{E}$
имеет по обеим сторонам заряженной поверхности различные значения
[уравнение (4.3)], и, стало быть, значение вектора $\mathbf{E}$ на самой
поверхности остаётся неопределённым.

В случае одного единственного уединённого проводника все электрические
силы сводятся ко взаимному отталкиванию элементов заряда этого
проводника. Ввиду того, что взаимно отталкивающиеся элементы заряда
не могут покинуть проводника, к поверхности проводника будут приложены
пондеромоторные силы, стремящиеся её растянуть\,\footnote{Если сообщить
электрический заряд мыльному пузырю, то под влиянием этих сил
отталкивания пузырь будет расширяться до тех пор, пока они не
уравновесятся силами поверхностного натяжения и разницей давления
воздуха внутри и вне пузыря.}. Подобного же рода силы будут, очевидно,
приложены и к поверхности неуединённого проводника, помещённого
в произвольное электростатическое поле. Чтобы определить величину этих
сил, рассмотрим некоторый элемент $dS$ поверхности проводника.
Напряжённость поля с \textit{внешней стороны} элемента $dS$, согласно
(7.8), равна:
$$
    E = 4\pi\sigma\mathbf{n}
$$
($\sigma$ — плотность его заряда), и направлена нормально к его
поверхности; внутри же проводника $E = 0$.

Однако $\mathbf{E}$ есть напряжённость результирующего поля всех
имеющихся зарядов, в том числе и заряда самого элемента $dS$:
$$
    \mathbf{E} = \mathbf{E}' + \mathbf{E}'',
$$
где $\mathbf{E}'$ есть поле элемента $dS$, а $\mathbf{E}''$ — поле
остальных зарядов. В двух смежных точках, лежащих по разным сторонам
элемента $dS$,

% стр. 90 (левая колонка)

поле $\mathbf{E}''$ этих зарядов будет, очевидно, одинаковым, поле же
$\mathbf{E}'$ будет иметь одинаковую величину $E'$, но различное
направление (рис. 22). Стало быть, с внешней стороны $E = E' + E'' =
4\pi\sigma$, а с внутренней $E = E' - E'' = 0$, откуда
$$
    E' = E'' = 2\pi\sigma.
$$

Сила, испытываемая зарядом $\sigma\,dS$ элемента $dS$, определяется
полем $\mathbf{E}''$ зарядов, лежащих вне этого элемента, и, стало
быть, равна:
$$
    E''\sigma\,dS = 2\pi\sigma^2\,dS = \frac{1}{2}E\sigma\,dS.
$$

Таким образом, на единицу поверхности проводника действует
пондеромоторная сила $f$:
\begin{equation}
    f = 2\pi\sigma^2 = \frac{1}{2}E\sigma = \frac{1}{8\pi}E^2,
    \tag{17.2}
\end{equation}
направленная, как легко убедиться, по внешней нормали к этой
поверхности. Эту величину $f$ можно, очевидно, назвать
\textit{поверхностной плотностью пондеромоторных сил}.

\textbf{3.} Рассмотрим ещё для полноты случай объёмного распределения
заряда с плотностью $\rho$, хотя в электростатическом поле в отсутствии
диэлектриков этот случай практически неосуществим\,\footnote{За
исключением неоднородных проводников, внутри которых действуют сторонние
электродвижущие силы (§\,38); на этом случае мы останавливаться не
будем.}. Если электрические заряды неразрывно связаны с элементами $dV$
объёма некоторого тела так, что перемещение заряда $\rho\,dV$ возможно
лишь при соответствующем перемещении элемента тела $dV$ и обратно, то
на каждый элемент этого тела будет действовать сила
\begin{equation}
    \mathbf{F} = \mathbf{f}\,dV = \mathbf{E}\rho\,dV,
    \tag{17.3}
\end{equation}
где $\mathbf{E}$ — напряжённость поля в элементе $dV$. Вычитать из
полного поля $\mathbf{E}$ поле заряда самого элемента $dV$ в этом случае
не нужно, ибо напряжённость поля бесконечно малого объёмного заряда
$\rho\,dV$ бесконечно мала даже внутри самого заряда и стремится к нулю
при беспредельном уменьшении его размеров. Проще всего убедиться в этом,
приняв, для простоты, что элемент $dV$ имеет форму шара, так что поле
его определяется формулой (4.8).

Из (17.3) следует, что \textit{объёмная плотность} пондеромоторных
сил равна\,\footnote{В дальнейшем плотность сил мы будем всегда
обозначать через малое $\mathbf{f}$, сохраняя $\mathbf{F}$ прописное
для обозначения сил, действующих на отдельный заряд, диполь, проводник
и т.\,д.}:
\begin{equation}
    \mathbf{f} = \rho\mathbf{E}.
    \tag{17.4}
\end{equation}

% стр. 91 (правая колонка) — Пример. Силы, действующие на диполь.

\medskip
\noindent\textbf{Пример.} \textit{Силы, действующие на диполь.}

Пусть $\mathbf{E}$ и $\mathbf{E}'$ суть напряжённости внешнего (по
отношению к диполю) поля в точках $P$ и $P'$, занимаемых соответственно
его отрицательным и положительным зарядами (рис. 23). Тогда
равнодействующая $\mathbf{F}$ сил поля, приложенных к этим зарядам,
будет равна:
$$
    \mathbf{F} = e\mathbf{E}' - e\mathbf{E} = e(\mathbf{E}' - \mathbf{E}).
$$
Разность $\mathbf{E}' - \mathbf{E}$ есть приращение вектора $\mathbf{E}$
на отрезке $PP'$, равном длине диполя $l$. Ввиду малости этого отрезка
приращение это может быть выражено, согласно (51*), формулой
$$
    \mathbf{E}' - \mathbf{E} = l\frac{\partial\mathbf{E}}{\partial l}
    = \mathbf{l}\nabla\cdot\mathbf{E},
$$
откуда
\begin{equation}
    \mathbf{F} = e\mathbf{l}\nabla\cdot\mathbf{E}
              = \mathbf{p}\nabla\cdot\mathbf{E}.
    \tag{17.5}
\end{equation}

Таким образом, величина силы $\mathbf{F}$, действующей в электрическом
поле на диполь, зависит от быстроты изменения этого поля в направлении
оси диполя; в однородном же поле силы, действующие на полюсы диполя,
равны по величине и противоположны по направлению и, стало быть, взаимно
уравновешиваются.

Помимо равнодействующей $\mathbf{F}$ для полного определения сил,
действующих на диполь, необходимо определить момент $\mathbf{N}$ этих
сил относительно центра диполя. Момент их относительно точки $P$,
в которой находится заряд $-e$, равен, очевидно:
$$
    \mathbf{N} = [\mathbf{l}\cdot e\mathbf{E}'] = [\mathbf{p}\mathbf{E}'].
$$
В пределе при достаточной малости $l$ точка $P$ совпадает с центром
диполя, а $\mathbf{E}'$ совпадает с $\mathbf{E}$, так что окончательно
\begin{equation}
    \mathbf{N} = [\mathbf{p}\mathbf{E}].
    \tag{17.6}
\end{equation}

Из этого выражения следует, что диполь стремится повернуться в
электрическом поле так, чтобы его момент $\mathbf{p}$ был параллелен
полю. При $\mathbf{p}$, направленном обратно $\mathbf{E}$, вращающий
момент также равен нулю, но это равновесие неустойчиво.

% ============================================================
% §18. ОПРЕДЕЛЕНИЕ ПОНДЕРОМОТОРНЫХ СИЛ ИЗ ВЫРАЖЕНИЯ ЭНЕРГИИ
% (стр. 91–95)
% ============================================================

\section*{\S\,18. Определение пондеромоторных сил из выражения энергии.}

% стр. 91 (правая колонка, продолжение)

\textbf{1.} Перемещение тел в электрическом поле сопровождается,
вообще говоря, как изменением энергии поля $W$ этого поля, так и
работой $A$ пондеромоторных сил поля. Если при этом не происходит
преобразования других форм энергии (тепловой, химической и т.\,д.), то

% стр. 92 (левая колонка)

работа пондеромоторных сил $A$ должна, очевидно, совершаться за счёт
изменения $\delta W$ энергии поля, так что
\begin{equation}
    A = -\delta W.
    \tag{18.1}
\end{equation}

Исходя из этого соотношения, мы и определили в §\,15 энергию поля
$W$ (15.1); обратно, зная $W$, мы на основании (18.1) можем определить
работу $A$, а стало быть, и величину совершающих эту работу
пондеромоторных сил.

Действительно, формула (18.1) показывает, что энергия электрического
поля $W$ играет роль \textit{потенциальной} энергии в смысле
аналитической механики, методами которой мы и можем воспользоваться.
Пусть энергия $W$ в самом общем случае выражена в зависимости от
каких-либо «обобщённых координат» $q_i$, характеризующих распределение
зарядов, проводников и т.\,п.

Тогда
\begin{equation}
    \delta W = \sum_i \frac{\partial W}{\partial q_i}\delta q_i
    \quad\text{и}\quad
    A = \sum_i \Theta_i\delta q_i,
    \tag{18.2}
\end{equation}
где $\Theta_i$, по терминологии аналитической механики, суть
«обобщённые силы», действующие «по направлению» координат $q_i$.
Отсюда на основании (18.1), при условии независимости координат $q_i$
следует, что
\begin{equation}
    \Theta_i = -\frac{\partial W}{\partial q_i}\,^*.
    \tag{18.3}
\end{equation}

Весьма часто определение пондеромоторных сил на основании этой
формулы оказывается несравненно более простым, чем непосредственное
определение их по формулам §\,17 путём интегрирования по отдельным
элементам зарядов.

\textbf{2.} В качестве простейшего примера определим этим способом
силы взаимного притяжения пластин плоского конденсатора. Энергия
этого конденсатора равна [см.\ уравнение (15.9) и результаты решения
задачи 8, стр. 55]:
$$
    W = \frac{e^2}{2C} = \frac{2\pi e^2 d}{S}.
$$

При раздвигании пластин, т.\,е. увеличении расстояния $d$ между
ними, совершается работа
$$
    A = -\delta W = -\frac{2\pi e^2}{S}\delta d.
$$

С другой стороны, если $f$ есть сила притяжения, испытываемая
единицей поверхности пластины конденсатора, а $F$ — общая сила,
действующая на всю поверхность пластины $S$, то
$$
    A = -F\cdot\delta d = -fS\cdot\delta d.
$$

% стр. 93 (правая колонка)

Приравнивая полученные выражения, найдём:
$$
    F = fS = \frac{2\pi e^2}{S} = 2\pi\sigma^2\cdot S,
$$
где $\sigma = \dfrac{e}{S}$ — плотность заряда на поверхности пластин,
что вполне согласуется с (17.2).

К тому же результату можно прийти и непосредственно из (18.3),
полагая в этой формуле $q = d$ и $\Theta = F$.

\textbf{3.} В предыдущем мы неявным образом предполагали, что при
смещении пластин конденсатора \textit{заряд} их $e$ остаётся
постоянным. Мы получили бы совершенно иное значение величины $\delta W$,
предположив, что не заряд, а \textit{разность потенциалов} пластин
конденсатора $\varphi_2 - \varphi_1$ остаётся при их смещении
постоянной. Объясняется это тем, что разность потенциалов обкладок
конденсатора может оставаться неизменной при перемещениях этих обкладок
(изменения ёмкости) лишь в том случае, если эта разность потенциалов
поддерживается неизменной некоторыми сторонними э.\,д.\,с.
неэлектростатического происхождения (см.\ §\,38) при изменении $C$
некоторую добавочную работу. (Это имеет место, например, при
присоединении обкладок конденсатора к полюсам гальванического
элемента.) При этом формула (18.1), очевидно, перестаёт быть применимой.

На основании (15.9) и (9.1) для конденсатора любой формы имеем:
\begin{align*}
    \text{при } e = \text{const}: \quad
    &\delta_e W = \delta\!\left(\frac{e^2}{2C}\right)
               = -\frac{e^2}{2C^2}\delta C
               = -\frac{1}{2}(\varphi_2 - \varphi_1)^2\delta C, \\
    \text{при } \varphi_2 - \varphi_1 = \text{const}: \quad
    &\delta_\varphi W = \delta\!\left\{\frac{1}{2}C(\varphi_2 - \varphi_1)^2\right\}
               = \frac{1}{2}(\varphi_2 - \varphi_1)^2\delta C,
\end{align*}
т.\,е.
\begin{equation}
    \delta_\varphi W = -\delta_e W = A.
    \tag{18.4}
\end{equation}

Таким образом, сопровождающие перемещение обкладок конденсатора
изменения его электрической энергии при $e = \text{const}$ и при
$\varphi_2 - \varphi_1 = \text{const}$ равны по величине, но
противоположны по знаку. С другой стороны, ввиду соотношения
$$
    e = (\varphi_2 - \varphi_1)C,
$$
изменение ёмкости $C$ при постоянном $\varphi_2 - \varphi_1$ должно
сопровождаться изменением заряда конденсатора $e$
$$
    \delta e = (\varphi_2 - \varphi_1)\delta C,
$$
т.\,е. перенесением заряда $\delta e$ с одной из обкладок конденсатора
на другую. При прохождении этого заряда $\delta e$ через включённый
между обкладками конденсатора гальванический элемент химическая энергия
этого элемента, как мы увидим в главе III, уменьшается, а э.\,д.\,с.
элемента $\mathcal{E}$ совершает работу
$$
    P = (\varphi_2 - \varphi_1)\delta e = (\varphi_2 - \varphi_1)^2\delta C,
$$
и, стало быть,
\begin{equation}
    P = 2\delta_\varphi W = A + \delta_\varphi W.
    \tag{18.5}
\end{equation}

% стр. 94 (левая колонка)

Таким образом, при $\varphi_2 - \varphi_1 = \text{const}$ работа $A$
пондеромоторных сил поля совершается не за счёт энергии поля $W$,
а за счёт химической энергии гальванического элемента (или энергии
другого источника сторонней э.\,д.\,с.). В частности, \textit{положительная}
работа сил поля $A$ сопровождается \textit{приращением} электрической
энергии $W$, происходящим также за счёт энергии гальванического элемента.

Таким образом, формула (18.1) справедлива лишь при условии, что
перемещение тел в электрическом поле не сопровождается работой
сторонних э.\,д.\,с. неэлектростатического происхождения. Кроме того,
условием применимости формулы (18.1) является, очевидно, \textit{достаточная
медленность перемещения тел}, а именно столь малая скорость перемещения,
чтобы в каждый данный момент процесса электрическое состояние системы
могло быть описано уравнениями электростатики (т.\,е. бесконечно мало
отличалось от равновесного).

\medskip
\noindent\textbf{Пример.} \textit{Определить силы, действующие на твёрдый
диполь, исходя из его энергии во внешнем поле.} Энергия диполя во
внешнем поле, согласно (15.8), равна
$$
    W = -\mathbf{p}\mathbf{E} = -pE\cos\vartheta
$$
и является функцией: а) координат центра диполя $x$, $y$, $z$ и б)
угла $\vartheta$ между осью диполя и заданным направлением электрического
поля $\mathbf{E}$ (численную величину момента диполя считаем
неизменной: «твёрдый» диполь). Согласно известным положениям
аналитической механики, «обобщённые силы», соответствующие этим
обобщённым координатам, представляют собой не что иное, как

а) равнодействующую приложенных к диполю сил:
\begin{equation}
    F_x = -\frac{\partial W}{\partial x},
    \quad
    F_y = -\frac{\partial W}{\partial y},
    \quad
    F_z = -\frac{\partial W}{\partial z},
    \tag{18.6}
\end{equation}

и б) момент приложенных к диполю сил:
\begin{equation}
    N = -\frac{\partial W}{\partial\vartheta},
    \tag{18.7}
\end{equation}
стремящийся увеличить угол $\vartheta$.

Из (15.8) и (18.6) следует, что
\begin{equation}
    \mathbf{F} = -\nabla W = \nabla(\mathbf{p}\mathbf{E}).
    \tag{18.8}
\end{equation}

Это выражение в случае \textit{электростатического} поля, которым
мы здесь только и ограничиваемся, лишь по своей форме отличается
от ранее выведенного выражения (17.5):
$$
    \mathbf{F} = \mathbf{p}\nabla\cdot\mathbf{E},
$$

% стр. 95 (правая колонка)

Действительно, для всякого постоянного (не зависящего от координат
$x$, $y$, $z$) вектора $\mathbf{p}$ имеет место соотношение
\begin{equation}
    \nabla(\mathbf{p}\mathbf{E}) = \mathbf{p}\nabla\cdot\mathbf{E}
    + [\mathbf{p}\operatorname{rot}\mathbf{E}].
    \tag{18.9}
\end{equation}
Чтобы доказать это соотношение, достаточно рассмотреть слагающую
$\nabla(\mathbf{p}\mathbf{E}) - \mathbf{p}\nabla\cdot\mathbf{E}$ хотя
бы по оси $x$:
\begin{align*}
    \frac{\partial}{\partial x}(\mathbf{p}\mathbf{E}) - \mathbf{p}\nabla\cdot E_x
    &= p_x\frac{\partial E_x}{\partial x}
     + p_y\frac{\partial E_y}{\partial x}
     + p_z\frac{\partial E_z}{\partial x}
     - \left(p_x\frac{\partial E_x}{\partial x}
           + p_y\frac{\partial E_x}{\partial y}
           + p_z\frac{\partial E_x}{\partial z}\right) = \\
    &= p_y\!\left(\frac{\partial E_y}{\partial x} - \frac{\partial E_x}{\partial y}\right)
     + p_z\!\left(\frac{\partial E_z}{\partial x} - \frac{\partial E_x}{\partial z}\right) = \\
    &= p_y\operatorname{rot}_z E - p_z\operatorname{rot}_y E
     = [\mathbf{p}\operatorname{rot}\mathbf{E}]_x.
\end{align*}

В интересующем нас случае электростатического поля вектор
$\mathbf{E} = -\operatorname{grad}\varphi$, и поэтому его ротор, как
ротор всякого градиента скаляра, равен нулю:
$\operatorname{rot}\mathbf{E} = 0$ [уравнение (7.6)].

Таким образом, в электростатическом поле действительно
\begin{equation}
    \nabla(\mathbf{p}\mathbf{E}) = \mathbf{p}\nabla\cdot\mathbf{E}
    \tag{18.10}
\end{equation}
и формула (18.8) эквивалентна формуле (17.5).

Обращаясь к моменту сил, приложенных к диполю, получаем из (15.8)
и (18.7):
$$
    N = -\frac{\partial W}{\partial\vartheta} = -pE\sin\vartheta.
$$
Ввиду того, что $\vartheta \leqslant \pi$, правая часть этого выражения
отрицательна, и, стало быть, согласно смыслу формулы (18.7), момент
$N$ стремится уменьшить угол $\vartheta$ между $\mathbf{p}$ и $\mathbf{E}$
(рис. 24). Следовательно, по величине и направлению вектор момента
сил определяется формулой
\begin{equation}
    \mathbf{N} = [\mathbf{p}\mathbf{E}],
    \tag{18.11}
\end{equation}
совпадающей с ранее выведенной нами формулой (17.6).



% 
% 
% 

\bigskip
\noindent\textbf{\S\,32. Пондеромоторные силы в диэлектриках.}

\medskip
\textbf{1.} На каждый элемент объёма диэлектрика в электрическом поле
должна, очевидно, действовать сила, равная сумме сил, приложенных к
отдельным молекулам диэлектрика. Заменим попрежнему эти молекулы
эквивалентными диполями и воспользуемся формулой не (17.5) а (18.9), определяющей
равнодействующую сил, приложенных к диполю:
$$
    \mathbf{F} = \mathbf{p}\nabla\cdot\mathbf{E}' \modmath{+ \left[\vec{p} \times rot \vec{E'}\right]},
$$
где $\mathbf{E}'$ означает напряжённость действующего на диполь поля.

\textit{Плотность сил}, приложенных к диэлектрику, т.\,е. сила,
отнесённая к единице объёма диэлектрика, будет, следовательно, равна:
\begin{equation}
    \mathbf{f} = \sum\mathbf{F} = \sum\mathbf{p}\nabla\cdot\mathbf{E}'
              = N\overline{\mathbf{p}\nabla\cdot\mathbf{E}'} \,\modmath{+ N\overline{\left[\vec{p} \times rot \vec{E'}\right]}},
    \tag{32.1}
\end{equation}
где суммирование должно быть произведено по всем диполям (молекулам),
находящимся в единице объёма; $N$ есть число молекул в единице объёма,
а черта означает среднее значение. В слабо поляризующихся диэлектриках
($a \ll 1$) можно в первом приближении, во-первых, \textit{среднее
значение произведения} $\mathbf{p}\nabla\cdot\mathbf{E}'$ заменить
\textit{произведением средних значений} $\mathbf{p}$ и
$\nabla\cdot\mathbf{E}'$ и, во-вторых, можно пренебречь

% стр. 150 (левая колонка)

разницей между $\overline{\nabla\cdot\mathbf{E}'}$ и $\nabla\cdot\mathbf{E}$,
где $\mathbf{E}$ есть напряжённость среднего макроскопического поля
(см.\ §\,28). В этом приближении
$$
    \mathbf{f} = N\mathbf{p}\nabla\cdot\mathbf{E} = \mathbf{P}\nabla\cdot\mathbf{E} \,\modmath{+ \left[\vec{P} \times rot \vec{E}\right]},
$$
или, на основании формулы (22.6),
$$
    \mathbf{f} = \frac{\varepsilon-1}{4\pi}\mathbf{E}\nabla\cdot\mathbf{E} \,\modmath{+ \left[\frac{\epsilon - 1}{4 \pi}\vec{E} \times rot \vec{E}\right]}.
$$
Далее, согласно (47*),
$$
    \mathbf{E}\nabla\cdot\mathbf{E}
    = \frac{1}{2}\nabla E^2 - [\mathbf{E}\operatorname{rot}\mathbf{E}].
$$
В электростатическом поле последний член равен нулю, так как
$\mathbf{E} = -\nabla\varphi$, а ротор градиента равен нулю [уравнение
$(42^*_1)$]. Таким образом, в электростатическом поле
\begin{equation}
    \mathbf{E}\nabla\cdot\mathbf{E} = \frac{1}{2}\cdot\nabla E^2,
    \tag{32.2}
\end{equation}
\begin{equation}
    \mathbf{f} = \frac{\varepsilon-1}{8\pi}\nabla E^2.
    \tag{32.3}
\end{equation}

Таким образом, в указанном приближении плотность пондеромоторных сил в
диэлектриках пропорциональна \textit{градиенту квадрата напряжённости
поля}; это и понятно, ибо, во-первых, в однородном поле сумма
приложенных к каждому диполю сил равна нулю, и во-вторых, по мере
возрастания поля возрастают не только силы поля, но и поляризация,
т.\,е. векторная сумма моментов диполей диэлектрика. Сила $\mathbf{f}$
направлена в сторону возрастания \textit{абсолютной величины} вектора
$\mathbf{E}$ \textit{независимо от направления этого вектора}; причина
этого лежит в том, что при изменении направления вектора $\mathbf{E}$
изменяется также и направление поляризации $\mathbf{P}$. Таким образом,
в электрическом поле диэлектрик \textit{увлекается в область наибольшей
напряжённости поля}. Этими пондеромоторными силами обусловливается,
например, притяжение заряжёнными проводниками кусочков бумаги,
бузиновых шариков и т.\,д.

\textbf{2.} Приведённый вывод формулы (32.3) основан на ряде
приближений и упрощений. Общее же выражение для пондеромоторных сил
может быть получено из энергии для электрического поля $W$ [формула
(30.1) или (30.4)] путём рассмотрения изменения этой энергии $\delta W$,
связанного с бесконечно малым произвольным (виртуальным) перемещением
$\mathbf{q}$ находящихся в поле тел (проводников и диэлектриков);
конечно, перемещение различных точек этих тел может быть различным, так
что $\mathbf{q}$ является произвольной, но непрерывной функцией точки.
При соблюдении условий, рассмотренных в §\,18, это

% стр. 151 (правая колонка)

изменение энергии должно быть равно взятой с обратным знаком работе $A$
пондеромоторных сил поля [уравнение (18.1)]:
$$
    \delta W = -A,
$$
которая в свою очередь очевидно, равна:
$$
    A = \int \mathbf{q}\mathbf{f}\,dV,
$$
где интеграл должен быть распространён по всему объёму \textit{полного
поля}, причём $\mathbf{f}\,dV$ есть сила, приложенная к элементу объёма
$dV$, а $\mathbf{q}$ — испытываемое им смещение. Следовательно,
\begin{equation}
    \delta W = -\int \mathbf{q}\mathbf{f}\,dV,
    \tag{32.4}
\end{equation}
где $\mathbf{f}$ — плотность пондеромоторных сил поля. Определив
$\delta W$ из выражения энергии, можно с помощью (32.4) найти
$\mathbf{f}$.

В этом состоит наиболее общий метод вычисления пондеромоторных сил,
которым мы неоднократно будем пользоваться в дальнейшем.

\textbf{3.} Применяя этот общий метод к вычислению пондеромоторных сил
электрического поля, мы в этом параграфе предположим, что как
диэлектрическая постоянная $\varepsilon$, так и напряжённость поля
повсюду непрерывны, т.\,е. что значение $\varepsilon$ плавно изменяется
в пограничных слоях, служащих разделом различных сред, и что свободные
заряды в пограничных слоях отсутствуют ($\sigma = 0$). Вопрос же о
силах, приложенных к поверхностям разрыва, будет рассмотрен в §\,34.

\redtext{На основании допущения об отсутствии поверхностей разрыва во всех
поверхностных интегралах, с которыми нам придётся встретиться в этом
параграфе, интегрирование будет распространяться лишь на внешнюю
граничную поверхность поля.} Кроме того, мы условимся рассматривать
\textit{полное поле}, так что все эти поверхностные интегралы обратятся
в нуль. \redtext{Вместе с тем обратятся в нуль и все интегралы вида}
$$
    \redmath{\int \operatorname{div}\mathbf{a}\,dV},
$$
так как, согласно теореме Гаусса (17*), интегралы эти могут быть
преобразованы в интегралы поверхностные:
\begin{equation}
    \redmath{\int \operatorname{div}\mathbf{a}\,dV = \oint a_n\,dS = 0}.
    \tag{32.5}
\end{equation}

\modtext{Для системы типа MenDrive мы не можем игнорировать силы,
приложенные к поверхностям разрыва. И поэтому мы не сможем
применять (32.5) в наших выкладках, отбрасывая, как это делал Тамм,
дивергентные интегралы, но будем вынуждены разбираться, каким компонентам
сил соответствуют отброшенные Таммом слагаемые}

\textbf{4.} Согласно (30.1) и (30.4) энергия электростатического поля,
ввиду предполагаемого нами в этом параграфе отсутствия поверхностных
зарядов ($\sigma = 0$), может быть выражена одним из следующих уравнений:
$$
    W_1 = \frac{1}{2}\int \varphi\rho\,dV
$$
и
$$
    W_2 = W_1 = \frac{1}{8\pi}\int \varepsilon E^2\,dV
              = \frac{1}{8\pi}\int \varepsilon(\operatorname{grad}\varphi)^2\,dV,
$$

% стр. 152 (левая колонка)

причём, конечно, $W_1 = W_2$. Следовательно, изменение энергии при
произвольном бесконечно малом перемещении $\mathbf{q}$ находящихся
в поле тел равно:
$$
    \delta W_1 = \frac{1}{2}\int \varphi\,\delta\rho\,dV
               + \frac{1}{2}\int \rho\,\delta\varphi\,dV
$$
или
$$
    \delta W_2 = \frac{1}{8\pi}\int (\operatorname{grad}\varphi)^2
                 \delta\varepsilon\,dV
               + \frac{1}{4\pi}\int \varepsilon\operatorname{grad}\varphi\cdot
                 \operatorname{grad}(\delta\varphi)\,dV,
$$
где $\delta\rho$, $\delta\varphi$, $\delta\varepsilon$, \ldots\ суть
изменения величин $\rho$, $\varphi$, $\varepsilon$, \ldots,
обусловленные перемещением $\mathbf{q}$.

Как известно, порядок выполнения операций дифференцирования и
варьирования не влияет на их результат, так что
$$
    \delta(\operatorname{grad}\varphi)
    = \delta\!\left(\mathbf{i}\frac{\partial\varphi}{\partial x}
                  + \mathbf{j}\frac{\partial\varphi}{\partial y}
                  + \mathbf{k}\frac{\partial\varphi}{\partial z}\right)
    = \mathbf{i}\frac{\partial(\delta\varphi)}{\partial x}
    + \mathbf{j}\frac{\partial(\delta\varphi)}{\partial y}
    + \mathbf{k}\frac{\partial(\delta\varphi)}{\partial z}
    = \operatorname{grad}(\delta\varphi).
$$
С другой стороны,
$$
    \varepsilon\operatorname{grad}\varphi = -\varepsilon\mathbf{E} = -\mathbf{D}.
$$
Следовательно, на основании $(43^*_a)$ подынтегральное выражение
последнего интеграла можно преобразовать следующим образом:
$$
    \varepsilon\operatorname{grad}\varphi\cdot\operatorname{grad}(\delta\varphi)
    = -\mathbf{D}\operatorname{grad}(\delta\varphi)
    = -\operatorname{div}(\mathbf{D}\,\delta\varphi)
    + \delta\varphi\operatorname{div}\mathbf{D}.
$$
\redtext{Интеграл первого слагаемого правой части, согласно (32.5), равен нулю},
так что, воспользовавшись уравнением (22.2), окончательно получаем:
$$
    \delta W_2 = \frac{1}{8\pi}\int E^2\delta\varepsilon\,dV
               + \int \rho\,\delta\varphi\,dV  \,\modmath{-\int \operatorname{div}(\mathbf{D}\,\delta\varphi)\,dV}.
$$
Ввиду того, что $\delta W_2$ равно $\delta W_1$, приращение энергии
$\delta W$ можно, очевидно, представить также в следующей форме:
$$
    \delta W = 2\delta W_1 - \delta W_2
    = \int \varphi\,\delta\rho\,dV
    - \frac{1}{8\pi}\int E^2\delta\varepsilon\,dV \,\modmath{+\int \operatorname{div}(\mathbf{D}\,\delta\varphi)\,dV}.
$$
Таким образом, вычисление $\delta W$ сведено нами к определению изменения
плотности свободного электричества $\rho$ и диэлектрической постоянной
$\varepsilon$ при виртуальном перемещении $\mathbf{q}$ находящихся
в поле тел.

% \modblock{

% В разделе 4 §32 (стр.~152, левая колонка) Тамм преобразует интеграл:

% \begin{equation}
% \frac{1}{4\pi}\int \varepsilon\operatorname{grad}\varphi\cdot
% \operatorname{grad}(\delta\varphi)\,dV
% \label{eq:tamm_integral_1}
% \end{equation}

% Используя тождество $\varepsilon\operatorname{grad}\varphi = -\mathbf{D}$ и
% формулу $(43^*_a)$ для дивергенции произведения:

% \begin{equation}
% -\mathbf{D}\operatorname{grad}(\delta\varphi) =
% -\operatorname{div}(\mathbf{D}\,\delta\varphi) +
% \delta\varphi\operatorname{div}\mathbf{D}
% \label{eq:divergence_identity_1}
% \end{equation}

% } % modblock

% \modblock{

% Тамм пишет:

% \begin{quote}
% \textit{«Интеграл первого слагаемого правой части, согласно (32.5), равен
% нулю»}
% \end{quote}

% и отбрасывает член:

% \begin{equation}
% \boxed{
% I_1 = -\int \operatorname{div}(\mathbf{D}\,\delta\varphi)\,dV
% }
% \label{eq:discarded_term_1}
% \end{equation}
% } % modblock

\modblock{
% \modtext{
\subsubsection{Физический смысл члена $I_1$}

Применяем теорему Гаусса с учётом внутренних границ (\ref{eq:gauss_discontinuous}):

\begin{equation}
I_1 = -\oint_{S_\infty} (\mathbf{D}\,\delta\varphi)_n\,dS
- \sum_j \int_{S_j} [\mathbf{D}\,\delta\varphi]_n\,dS_j
\label{eq:term_1_gauss}
\end{equation}

Для MenDrive с границами $x = \pm A$:

\begin{equation}
I_1 = -\int_{x=-A} [D_x\,\delta\varphi]\,dS - \int_{x=+A} [D_x\,\delta\varphi]\,dS
\label{eq:term_1_mendrive}
\end{equation}

Вариация потенциала при виртуальном перемещении $\mathbf{q} = q_x\hat{x}$:

\begin{equation}
\delta\varphi = \varphi(\vec{r}+\mathbf{q}) - \varphi(\vec{r}) \approx
\mathbf{q}\cdot\nabla\varphi = -q_x E_x
\label{eq:potential_variation}
\end{equation}

На границе раздела сред $D_x$ \textit{непрерывна} (нет свободных зарядов),
но $E_x$ \textit{терпит разрыв}:

\begin{equation}
D_x^{(+)} = D_x^{(-)} = D_x, \quad
E_x^{(+)} = \frac{D_x}{\varepsilon_+}, \quad
E_x^{(-)} = \frac{D_x}{\varepsilon_-}
\label{eq:boundary_conditions_D_E}
\end{equation}

Скачок произведения:

\begin{equation}
\begin{aligned}
[D_x\,\delta\varphi] &= D_x(-q_x E_x^{(+)}) - D_x(-q_x E_x^{(-)}) \\
&= q_x D_x (E_x^{(-)} - E_x^{(+)}) \\
&= q_x D_x^2 \left(\frac{1}{\varepsilon_-} - \frac{1}{\varepsilon_+}\right) \\
&= q_x \frac{D_x^2(\varepsilon_+ - \varepsilon_-)}{\varepsilon_+\varepsilon_-}
\end{aligned}
\label{eq:jump_D_delta_phi}
\end{equation}

\subsubsection{Сила, соответствующая члену $I_1$}

Из $\delta W = -\int\mathbf{q}\cdot\mathbf{f}\,dV$ получаем силу на единицу
площади границы:

\begin{equation}
\boxed{
f_x^{(1)} = \frac{D_x^2}{8\pi}\left(\frac{1}{\varepsilon_-} - \frac{1}{\varepsilon_+}\right)
= \frac{\varepsilon_+ - \varepsilon_-}{8\pi\varepsilon_+\varepsilon_-}D_x^2
}
\label{eq:force_term_1}
\end{equation}

Для границы вакуум/диэлектрик ($\varepsilon_- = 1$, $\varepsilon_+ = \varepsilon$):

\begin{equation}
f_x^{(1)} = \frac{\varepsilon - 1}{8\pi\varepsilon}D_x^2 =
\frac{\varepsilon - 1}{8\pi}\varepsilon E_x^2
\label{eq:force_term_1_vacuum_dielectric}
\end{equation}

% modtext
% modblock
}


\subsubsection{Сравнение с конвективной силой на поляризацию}
\modblock{
Конвективная сила на поляризацию из формулы (18.9) Тамма:

\begin{equation}
f_{P,\text{conv}} = (\vec{P}\cdot\nabla)E_x =
\frac{\varepsilon-1}{4\pi}E_x\frac{\partial E_x}{\partial x}
\label{eq:convective_force_polarization}
\end{equation}

Интеграл по переходному слою от $x = A^-$ до $x = A^+$:

\begin{equation}
\begin{aligned}
F_{P,\text{conv}} &= \int_{A^-}^{A^+} \frac{\varepsilon-1}{4\pi}E_x\frac{\partial E_x}{\partial x}\,dx \\
&= \frac{\varepsilon-1}{8\pi}\left[E_x^2\right]_{A^-}^{A^+} \\
&= \frac{\varepsilon-1}{8\pi}\left(E_x^{(+)2} - E_x^{(-)2}\right)
\end{aligned}
\label{eq:convective_integrated}
\end{equation}

Используя $D_x = \varepsilon E_x^{(+)} = E_x^{(-)}$ (для вакуума $\varepsilon=1$):

\begin{equation}
F_{P,\text{conv}} = \frac{\varepsilon-1}{8\pi}\left(\frac{D_x^2}{\varepsilon^2} - D_x^2\right) =
-\frac{\varepsilon-1}{8\pi\varepsilon}D_x^2
\label{eq:convective_via_D}
\end{equation}

\textbf{Вывод:} Член $I_1$, отброшенный Таммом, в точности соответствует
\textit{конвективной силе на связанные заряды поляризации} на границе
раздела сред! Для MenDrive с ферритом это даёт вклад в тягу
($\sim $ Н/кВт для ветви 5).
}



\textbf{5.} Величины $\delta\rho$ и $\delta\varepsilon$, входящие в
последнее уравнение, носят название \textit{локальных изменений} величин
$\rho$ и $\varepsilon$, в отличие от так называемых \textit{изменений
материальных}, которые мы обозначим через $\delta'\rho$ и $\delta'\varepsilon$.
Разница между этими понятиями состоит в следующем. \textit{Локальное}
изменение, испытываемое какой-либо величиной при перемещениях среды,
есть изменение значения этой величины в некоторой определённой точке
пространства, не принимающей участия в этом перемещении.
\textit{Материальное} же изменение есть изменение значения данной
величины в некотором определённом материальном элементе перемещающейся
среды. Иными словами, если элемент $dV$ среды, находившийся в точке $P$,
перемещается в точку $P'$, то $\delta\rho$ равно разности между
значениями $\rho$ в одной и той же точке $P$ до и после перемещения,
а $\delta'\rho$ равно разности между значениями, которыми обладает $\rho$
в точке $P'$ после перемещения и в точке $P$ до перемещения.

Соответственно этому локальное изменение $\delta\rho$ какой-либо
величины в точке $P$ может быть представлено в виде суммы двух слагаемых.
Во-первых, после смещения $\mathbf{q}$ в точке $P$ будет находиться
элемент объёма среды, ранее

% стр. 153 (правая колонка)

находившийся в точке $P''$ на расстоянии $-\mathbf{q}$ от $P$, в которой
величина $\rho$ имела, очевидно, значение
$$
    \rho_{P''} = \rho_P + (-\mathbf{q})\operatorname{grad}\rho
              = \rho_P - \mathbf{q}\operatorname{grad}\rho.
$$
Кроме того, значение величины $\rho$ в этом объёме испытает материальное
изменение $\delta'\rho$. Следовательно\,\footnote{Ввиду предполагаемой
бесконечности малости смещения $\mathbf{q}$ мы вправе не учитывать
разницы между $(\delta'\rho)_P$ и $(\delta'\rho)_{P'}$, сводящейся к
величинам второго порядка малости.},
\begin{equation}
    \delta\rho = \delta'\rho - \mathbf{q}\operatorname{grad}\rho.
    \tag{32.6}
\end{equation}
Таким образом, определение локального изменения сведено к определению
изменения материального.

Если перемещение среды не сопровождается её деформацией, то
материальное изменение $\delta'\rho$ равно, очевидно, нулю. Если же
элемент $V$ среды, обладающий зарядом $de$, испытывает при смещении
сжатие или расширение, так что объём его становится равным $V + \delta'V$,
то, очевидно,
$$
    de = \rho V = (\rho + \delta'\rho)(V + \delta'V).
$$
Произведя приведение членов и пренебрегая произведением
$\delta'\rho\cdot\delta'V$, получаем:
$$
    \delta'\rho = -\rho\frac{\delta'V}{V}.
$$
Чтобы определить относительное изменение $\delta'V/V$ объёма элемента
среды, заметим, что при смещении $\mathbf{q}$ элемента $dS$ поверхности
$S$, ограничивающей объём $V$, величина этого объёма возрастает на
$q_n\,dS$, т.\,е. на величину цилиндра, описанного элементом $dS$ при
этом перемещении. Конечно, $q_n\,dS$ может иметь как положительное, так
и отрицательное значение в зависимости от величины угла между
$\mathbf{q}$ и внешней нормалью $\mathbf{n}$ к элементу $dS$. Полное же
изменение объёма $V$ при смещении $\mathbf{q}$ будет, очевидно, равно:
$$
    \delta'V = \oint q_n\,dS,
$$
и, следовательно [см.\ определение дивергенции вектора, уравнение (18*)],
в пределе при $V \to 0$:
\begin{equation}
    \frac{\delta'V}{V} = \frac{1}{V}\oint q_n\,dS
                       = \operatorname{div}\mathbf{q}.
    \tag{32.7}
\end{equation}
Внося это значение в выражение для $\delta'\rho$, получаем:
\begin{equation}
    \delta'\rho = -\rho\operatorname{div}\mathbf{q}.
    \tag{32.8}
\end{equation}
Очевидно, что формулой этого вида будет определяться не только
локальное изменение плотности электричества $\rho$, но, например, и
изменение плотности $\tau$ самой среды (т.\,е. изменение массы единицы
объёма):
$$
    \delta'\tau = -\tau\operatorname{div}\mathbf{q}.
$$

\textbf{6.} Внося (32.8) в (32.6), получаем с помощью $(43^*_a)$
\begin{equation}
    \delta\rho = -\rho\operatorname{div}\mathbf{q}
               - \mathbf{q}\operatorname{grad}\rho
               = -\operatorname{div}(\rho\mathbf{q}).
    \tag{32.9}
\end{equation}
С другой стороны, по аналогии с (32.6), имеем:
$$
    \delta\varepsilon = \delta'\varepsilon - \mathbf{q}\operatorname{grad}\varepsilon.
$$

% стр. 154 (левая колонка)

Так как значение $\varepsilon$ зависит от плотности
диэлектрика\,\footnote{В твёрдых диэлектриках $\varepsilon$ может
зависеть (при заданной температуре) не только от плотности и не только
от его деформаций, не связанных с изменениями плотности. Исследование
этого обстоятельства требует учёта анизотропии деформированного
диэлектрика и проведено, например, в книге
С\,т\,р\,е\,т\,т\,о\,н\,а, Теория электромагнетизма, Гостехиздат,
1948\,г.}, то материальное изменение $\varepsilon$ будет определяться
материальным изменением $\delta'\tau$ плотности $\tau$:
$$
    \delta'\varepsilon
    = \frac{\partial\varepsilon}{\partial\tau}\delta'\tau
    = -\frac{\partial\varepsilon}{\partial\tau}\tau\operatorname{div}\mathbf{q}
$$
и, стало быть,
\begin{equation}
    \delta\varepsilon
    = -\frac{\partial\varepsilon}{\partial\tau}\tau\operatorname{div}\mathbf{q}
    - \mathbf{q}\operatorname{grad}\varepsilon.
    \tag{32.9$'$}
\end{equation}
Внося значения $\delta\varepsilon$ и $\delta\rho$ [уравнения (32.9)]
в выражение для $\delta W$, получим:
$$
    \delta W = \frac{1}{8\pi}\int E^2\operatorname{grad}\varepsilon\cdot\mathbf{q}\,dV
             + \frac{1}{8\pi}\int E^2\frac{\partial\varepsilon}{\partial\tau}
               \tau\operatorname{div}\mathbf{q}\,dV
             - \int \varphi\operatorname{div}(\rho\mathbf{q})\,dV.
$$
Подынтегральные выражения двух последних интегралов можно преобразовать
с помощью $(43^*_a)$:
$$
    E^2\frac{\partial\varepsilon}{\partial\tau}\tau\operatorname{div}\mathbf{q}
    = \operatorname{div}\!\left(E^2\frac{\partial\varepsilon}{\partial\tau}
      \tau\mathbf{q}\right)
    - \mathbf{q}\operatorname{grad}\!\left(E^2\frac{\partial\varepsilon}{\partial\tau}
      \tau\right)
$$
и
$$
    \varphi\operatorname{div}(\rho\mathbf{q})
    = \operatorname{div}(\varphi\rho\mathbf{q})
    - \rho\mathbf{q}\operatorname{grad}\varphi.
$$
Согласно (32.5) интегралы первых слагаемых этих выражений обращаются
в нуль; приняв ещё во внимание, что $\operatorname{grad}\varphi = -\mathbf{E}$,
получаем:
$$
    \delta W = \int \mathbf{q}\left\{
        \frac{1}{8\pi}E^2\operatorname{grad}\varepsilon
        - \frac{1}{8\pi}\operatorname{grad}\!\left(
            E^2\frac{\partial\varepsilon}{\partial\tau}\tau
          \right)
        - \rho\mathbf{E}
    \right\}dV.
$$

\textbf{7.} Это выражение для $\delta W$ должно совпадать с выражением
(32.4) при произвольной зависимости смещения $\mathbf{q}$ от координат
точки. Отсюда следует равенство выражений, на которые множится
$\mathbf{q}$ в обоих интегралах:
\begin{equation}
    \mathbf{f} = \rho\mathbf{E}
               - \frac{1}{8\pi}E^2\operatorname{grad}\varepsilon
               + \frac{1}{8\pi}\operatorname{grad}\!\left(
                   E^2\frac{\partial\varepsilon}{\partial\tau}\tau
                 \right).
    \tag{32.10}
\end{equation}
Эта формула и является искомым выражением для плотности пондеромоторных
сил $\mathbf{f}$. Она слагается из двух частей $\mathbf{f}^{(1)}$
и $\mathbf{f}^{(2)}$, а именно из $\mathbf{f}^{(1)}$:
\begin{equation}
    \mathbf{f}^{(1)} = \rho\mathbf{E},
    \tag{32.11}
\end{equation}
действующей на свободные электрические заряды, и из
\begin{equation}
    \mathbf{f}^{(2)}
    = \frac{1}{8\pi}\operatorname{grad}\!\left(
        E^2\frac{\partial\varepsilon}{\partial\tau}\tau
      \right)
    - \frac{1}{8\pi}E^2\operatorname{grad}\varepsilon,
    \tag{32.12}
\end{equation}
зависящей от $\frac{\partial\varepsilon}{\partial\tau}$ и
$\operatorname{grad}\varepsilon$ и отличной от нуля только в диэлектриках.

% стр. 155 (правая колонка)

Из (32.11) следует, что плотность $\mathbf{f}^{(1)}$ пондеромоторных
сил, действующих на \textit{свободные} заряды, в диэлектрике, как и
в вакууме, определяется напряжённостью электрического поля. Из этого
положения мы исходили в §\,30 при подсчёте работы, совершаемой полем
при перемещении свободных зарядов. Что же касается плотности
$\mathbf{f}^{(2)}$ пондеромоторных сил, действующих на диэлектрик, то
в слабо поляризующихся диэлектриках ($a \ll 1$) выражение (32.12)
совпадает с тем, которое было получено нами непосредственным подсчётом
пондеромоторных сил в этих диэлектриках.

Действительно, согласно (29.12),
$$
    \frac{\varepsilon-1}{\varepsilon+2} = \frac{C\tau}{3},
$$
где коэффициент $C$ от плотности диэлектрика $\tau$ не зависит. При
$a \ll 1$, т.\,е. при $\varepsilon \approx 1$, можно с достаточной
точностью положить
$$
    \varepsilon - 1 = C\tau.
$$
Тогда
$$
    \tau\frac{\partial\varepsilon}{\partial\tau} = \tau C = \varepsilon - 1,
$$
и (32.12) принимает вид:
$$
    \mathbf{f}^{(2)} = \frac{1}{8\pi}\operatorname{grad}\{E^2(\varepsilon-1)\}
                     - \frac{1}{8\pi}E^2\operatorname{grad}\varepsilon.
$$
Но, согласно $(43^*_a)$,
$$
    \operatorname{grad}\{E^2(\varepsilon-1)\}
    = (\varepsilon-1)\operatorname{grad}E^2 + E^2\operatorname{grad}\varepsilon.
$$
Следовательно,
$$
    \mathbf{f}^{(2)} = \frac{\varepsilon-1}{8\pi}\operatorname{grad}E^2,
$$
что, действительно, совпадает с (32.3).

Таким образом, условием применимости формулы (32.3) является линейная
зависимость диэлектрической постоянной от плотности диэлектрика,
имеющая место, строго говоря, только в газах.

Заметим в заключение, что Максвелл и ряд других авторов, например
Абрагам, не принимали во внимание зависимости диэлектрической постоянной
от плотности среды, благодаря чему выражение пондеромоторных сил
в диэлектриках, которым они пользовались:
\begin{equation}
    \mathbf{f} = -\frac{1}{8\pi}E^2\operatorname{grad}\varepsilon,
    \tag{32.13}
\end{equation}
отличалось от (32.12) отсутствием первого члена
\begin{equation}
    \mathbf{f}'' = \frac{1}{8\pi}\operatorname{grad}\!\left(
        E^2\frac{\partial\varepsilon}{\partial\tau}\tau
    \right).
    \tag{32.14}
\end{equation}

% стр. 156 (левая колонка)

Этот член $\mathbf{f}''$ в твёрдых и жидких диэлектриках сравним по
величине с (32.13), так что пренебрегать им, вообще говоря, не
представляется возможным.

Однако надо иметь в виду, что отличие формулы (32.12) от максвелловой
формулы (32.13) сказывается лишь на распределении сил по объёму
диэлектрика; результирующая же не учтённых Максвеллом сил $\mathbf{f}''$,
приложенных к какому-либо телу, либо равна нулю (если это тело помещено
в вакуум), либо уравновешивается гидростатическим давлением, возникающим
в окружающей среде под влиянием электрического поля. Это утверждение
будет доказано в §\,34.

\section*{\S\,33. Сведение объёмных сил к натяжениям\,\footnote{Этот
и следующий параграфы можно опустить при первом чтении книги.}.}

\textbf{1.} Как уже упоминалось в §\,16, механистическая теория
электромагнитного поля прошлого века искала причины электрических
явлений в упругих деформациях гипотетической среды — эфира.
Характерной особенностью сил упругости, как, впрочем, и вообще сил
близкодействия, является возможность сведения их к \textit{натяжениям},
т.\,е. возможность сведения сил, действующих на произвольный участок
среды, к силам натяжения, испытываемым \textit{поверхностью} этого
участка (в частности, давление есть отрицательное натяжение).

Соответственно этому перед механистической теорией поля стояла задача
сведения пондеромоторных сил поля к упругим натяжениям среды. Свести
эти силы к натяжениям, как мы покажем, действительно, оказывается
возможным. Правда, это обстоятельство ни в коей мере не спасает
механистической теории поля, оказавшейся в целом несостоятельной;
однако при рассмотрении многих вопросов замена пондеромоторных сил
эквивалентными им натяжениями оказывается весьма целесообразной.

В этом параграфе мы рассмотрим вопрос о сведении объёмных сил
к натяжениям в общей форме с тем, чтобы в следующем параграфе применить
полученные результаты к интересующему нас случаю электрического поля.

\textbf{2.} Рассмотрим некоторый объём среды $V$, ограниченный
поверхностью $S$. Если $\mathbf{f}$ есть объёмная плотность сил, то
равнодействующая всех сил, приложенных к телам, находящимся внутри
объёма $V$, будет равна:
\begin{equation}
    \mathbf{F} = \int_V \mathbf{f}\,dV.
    \tag{33.1}
\end{equation}
С другой стороны, если объёмные силы вообще могут быть сведены

% стр. 157 (правая колонка)

к натяжениям, то той же величине должна равняться и совокупность
натяжений, действующих \textit{извне} на замкнутую поверхность $S$.

Сила натяжения, испытываемая каким-либо элементом поверхности
произвольного участка среды, пропорциональна величине этого элемента
$dS$ и зависит не только от положения этого элемента, но и от его
направления нормали $\mathbf{n}$ к элементу $dS$. В частности, при
повороте площадки $dS$ на $180^\circ$, т.е. при изменении направления
нормали $\mathbf{n}$ на обратное, сила натяжения меняет свой знак.
В этом выражается тот факт, что на противоположные стороны
произвольного элемента поверхности $dS$\,\footnote{За исключением
элементов поверхностей разрыва, которых мы в этом параграфе не
рассматриваем.} всегда действуют взаимно уравновешивающиеся равные
и противоположные натяжения.

Обозначим через $\mathbf{T}_n$ силу, действующую извне на единицу
поверхности, внешняя нормаль к которой направлена по $\mathbf{n}$;
компоненты этой силы по осям координат будем обозначать через $T_{xn}$,
$T_{yn}$ и $T_{zn}$\,\footnote{Часто эти компоненты обозначают через
$T_{nx}$, $T_{ny}$, $T_{nz}$. Произведённый нами выбор порядка
индексов соответствует принятому в общей теории, изложенной в
§\,105.}. Тогда равнодействующая всех натяжений, приложенных извне
к поверхности, выразится, очевидно, через
\begin{equation}
    \mathbf{F} = \oint_S \mathbf{T}_n\,dS,
    \tag{33.2}
\end{equation}
где $\mathbf{n}$ есть внешняя нормаль к элементу $dS$. Приравнивая
выражения (33.1) и (33.2)
\begin{equation}
    \mathbf{F} = \int_V \mathbf{f}\,dV = \oint_S \mathbf{T}_n\,dS,
    \tag{33.3}
\end{equation}
мы можем найти соотношение между плотностью объёмных сил $\mathbf{f}$
и натяжений $\mathbf{T}_n$.

\textbf{3.} Выберем какую-либо произвольную систему декартовых координат
и обозначим соответственно через $\mathbf{T}_x$, $\mathbf{T}_y$ и
$\mathbf{T}_z$ векторную величину силы натяжения, действующей извне на
единичную площадку, внешняя нормаль к которой направлена соответственно
по оси $x$, по оси $y$ и по оси $z$. Слагающие этих сил по осям
координат обозначим через $T_{xx}$, $T_{yx}$, $T_{zx}$ и т.\,д., так
что
\begin{align*}
    \mathbf{T}_x &= \mathbf{i}T_{xx} + \mathbf{j}T_{yx} + \mathbf{k}T_{zx}, \\
    \mathbf{T}_y &= \mathbf{i}T_{xy} + \mathbf{j}T_{yy} + \mathbf{k}T_{zy}, \\
    \mathbf{T}_z &= \mathbf{i}T_{xz} + \mathbf{j}T_{yz} + \mathbf{k}T_{zz}.
\end{align*}
Таким образом, например, $T_{zz}$ есть слагающая по оси $x$ силы
$\mathbf{T}_z$, действующей на единичную площадку, внешняя нормаль
к которой направлена по оси $z$.

% стр. 158 (левая колонка)

Легко показать (см.\ любой учебник теории упругости), что сила
натяжения $\mathbf{T}_n$, действующая извне на произвольно
ориентированную единичную площадку, внешняя нормаль к которой имеет
направление $n$, связана с силами $\mathbf{T}_x$, $\mathbf{T}_y$
и $\mathbf{T}_z$ следующим соотношением:
\begin{equation}
    \mathbf{T}_n = \mathbf{T}_x\cos(n,x) + \mathbf{T}_y\cos(n,y)
                 + \mathbf{T}_z\cos(n,z),
    \tag{33.4}
\end{equation}
так что слагающая этой силы, например по оси $x$, равна:
\begin{equation}
    T_{xn} = T_{xx}\cos(n,x) + T_{xy}\cos(n,y) + T_{xz}\cos(n,z).
    \tag{33.5}
\end{equation}

Таким образом, девять величин $T_{xx}$, $T_{xy}$ и т.\,д. вполне
характеризуют собою систему натяжений в данной точке пространства:
значениями этих величин вполне определяются натяжения, испытываемые
произвольно ориентированной площадкой. Величины $T_{xx}$, $T_{xy}$
и т.\,д. называются \textit{компонентами тензора натяжений}, а их
совокупность носит название \textit{тензора натяжений}, который будет
обозначаться нами буквой $\mathbf{T}$ (без индексов). Компоненты
тензора могут быть записаны в следующей симметричной форме:
\begin{equation}
    \mathbf{T} =
    \begin{pmatrix}
        T_{xx} & T_{xy} & T_{xz} \\
        T_{yx} & T_{yy} & T_{yz} \\
        T_{zx} & T_{zy} & T_{zz}
    \end{pmatrix}.
    \tag{33.6}
\end{equation}

\textbf{4.} Чтобы от интегрального соотношения (33.3) между натяжениями
и объёмными силами перейти к соотношениям дифференциальным, мы должны,
очевидно, прежде всего преобразовать поверхностный интеграл справа
в объёмный (или наоборот). Воспользовавшись формулами (33.2) и (33.5),
мы можем следующим образом выразить слагающую равнодействующей
$\mathbf{F}$ по оси $x$ через компоненты тензора натяжений:
$$
    F_x = \oint_S T_{xn}\,dS
        = \oint_S \{T_{xx}\cos(n,x) + T_{xy}\cos(n,y) + T_{xz}\cos(n,z)\}\,dS.
$$
Воспользуемся теперь теоремой Гаусса (17*), которая в развёрнутой
форме гласит:
$$
    \oint_S \{a_x\cos(n,x) + a_y\cos(n,y) + a_z\cos(n,z)\}\,dS
    = \int_V\!\left(
        \frac{\partial a_x}{\partial x}
      + \frac{\partial a_y}{\partial y}
      + \frac{\partial a_z}{\partial z}
    \right)dV.
$$
Так как эта теорема справедлива для любых непрерывных функций точки
$a_x$, $a_y$, $a_z$, то, полагая в ней $a_x = T_{xx}$, $a_y = T_{xy}$,
$a_z = T_{xz}$ (при условии непрерывности компонент тензора натяжений
внутри объёма $V$):
$$
    F_x = \oint_S T_{xn}\,dS
        = \int_V\!\left(
            \frac{\partial T_{xx}}{\partial x}
          + \frac{\partial T_{xy}}{\partial y}
          + \frac{\partial T_{xz}}{\partial z}
          \right)dV.
$$

% стр. 159 (правая колонка)

Внося это в (33.3) и приравнивая ввиду произвольности объёма $V$
подынтегральные выражения, получаем окончательно:
\begin{equation}
    f_x = \frac{\partial T_{xx}}{\partial x}
        + \frac{\partial T_{xy}}{\partial y}
        + \frac{\partial T_{xz}}{\partial z},
    \quad
    f_y = \frac{\partial T_{yx}}{\partial x}
        + \frac{\partial T_{yy}}{\partial y}
        + \frac{\partial T_{yz}}{\partial z},
    \quad
    f_z = \frac{\partial T_{zx}}{\partial x}
        + \frac{\partial T_{zy}}{\partial y}
        + \frac{\partial T_{zz}}{\partial z}.
    \tag{33.7}
\end{equation}
Эти формулы и устанавливают искомые дифференциальные соотношения между
плотностью объёмных сил $\mathbf{f}$ и компонентами тензора натяжений
$\mathbf{T}$.

\textbf{5.} Из (33.7) следует, что плотность объёмных сил определяется
не абсолютной величиной натяжений, а характером изменения натяжений
в пространстве (при перемещениях точки наблюдения). В частности,
$\mathbf{f}$ равно нулю, если компоненты тензора натяжения $\mathbf{T}$
имеют в данном участке среды постоянные значения. Это и понятно, ибо
если мы мысленно выделим в среде произвольный параллелепипед, то в
случае постоянства тензора $\mathbf{T}$ на противоположных его гранях
будут действовать натяжения, равные по величине и противоположные по
направлению; следовательно, равнодействующая приложенных к
параллелепипеду сил будет равна нулю.

\textbf{6.} Для эквивалентности объёмных сил и натяжений необходимо,
чтобы при замене объёмных сил эквивалентными натяжениями оставались
неизменными не только \textit{равнодействующая} сил, приложенных к
произвольному объёму, но и \textit{момент} этих сил. Это обстоятельство
накладывает дополнительное ограничение на компоненты тензора натяжений.

Момент $\mathbf{N}$ объёмных сил, приложенных к произвольному объёму
$V$, равен
\begin{equation}
    \mathbf{N} = \int_V [\mathbf{R}\mathbf{f}]\,dV,
    \tag{33.8}
\end{equation}
где $\mathbf{R}$ есть расстояние от точки $O$, относительно которой
определяется момент сил, до элемента $dV$. Если объёмные силы
$\mathbf{f}$ эквивалентны натяжениям $\mathbf{T}$, то должно
выполняться соотношение (33.7). Следовательно, слагающая $N_x$,
например по оси $x$, должна равняться:
$$
    N_x = \int_V (yf_z - zf_y)\,dV =
$$
$$
    = \int_V\!\left\{
        y\!\left(
            \frac{\partial T_{zx}}{\partial x}
          + \frac{\partial T_{zy}}{\partial y}
          + \frac{\partial T_{zz}}{\partial z}
        \right)
        - z\!\left(
            \frac{\partial T_{yx}}{\partial x}
          + \frac{\partial T_{yy}}{\partial y}
          + \frac{\partial T_{yz}}{\partial z}
        \right)
    \right\}dV.
$$

% стр. 160 (левая колонка)

Подынтегральное выражение справа может быть представлено следующим
образом:
$$
    \frac{\partial}{\partial x}(yT_{zx} - zT_{yx})
  + \frac{\partial}{\partial y}(yT_{zy} - zT_{yy})
  + \frac{\partial}{\partial z}(yT_{zz} - zT_{yz})
  - T_{zy} + T_{yz}.
$$
Так как три первых члена этого выражения по своему виду совпадают
с выражением дивергенции вектора с компонентами
$$
    a_x = yT_{zx} - zT_{yx},\quad
    a_y = yT_{zy} - zT_{yy},\quad
    a_z = yT_{zz} - zT_{yz},
$$
то объёмный интеграл можно преобразовать с помощью теоремы Гаусса (17*):
$$
    N_x = \oint_S a_n\,dS + \int_V (T_{yz} - T_{zy})\,dV,
$$
причём на основании (33.5)
\begin{align*}
    a_n &= a_x\cos(x,n) + a_y\cos(y,n) + a_z\cos(z,n) = \\
        &= y\{T_{zx}\cos(x,n) + T_{zy}\cos(y,n) + T_{zz}\cos(z,n)\} - \\
        &\quad - z\{T_{yx}\cos(x,n) + T_{yy}\cos(y,n) + T_{yz}\cos(z,n)\} = \\
        &= yT_{zn} - zT_{yn}.
\end{align*}
Окончательно получаем:
\begin{equation}
    N_x = \int_V (yf_z - zf_y)\,dV
        = \oint_S (yT_{zn} - zT_{yn})\,dS
        + \int_V (T_{yz} - T_{zy})\,dV.
    \tag{33.9}
\end{equation}

Поверхностный интеграл справа равен моменту сил натяжения $\mathbf{T}_n$,
приложенных к поверхности $S$ объёма $V$. Момент этих сил натяжения
будет равняться моменту сил объёмных в том и только том случае, если
последний интеграл равен нулю. Ввиду произвольности объёма $V$ это
будет иметь место только при равенстве нулю подынтегрального выражения
во всех точках пространства:
$$
    T_{yz} = T_{zy}.
$$
Повторив те же рассуждения для слагающих $\mathbf{N}$ по осям $y$ и $z$,
получим следующие соотношения:
\begin{equation}
    T_{yz} = T_{zy},\quad T_{zx} = T_{xz},\quad T_{xy} = T_{yx}.
    \tag{33.10}
\end{equation}
Тензоры, компоненты которых удовлетворяют соотношениям (33.10),
называются \textit{симметричными}.

% стр. 161 (правая колонка)

натяжений не может быть заменена эквивалентным распределением
объёмных сил.

Это, впрочем, явствует уже из того, что если компоненты тензора
$\mathbf{T}$ постоянны по величине, то объёмные силы, согласно (33.7),
обращаются в нуль, тогда как момент сил натяжения, приложенных
к произвольному объёму, будет при $T_{lk} \neq T_{kl}$ отличаться
от нуля даже при постоянстве $T_{lk}$\,\footnote{Заметим, что
в анизотропных средах тензор натяжений электрического поля, вообще
говоря, не симметричен.}.

\textbf{7.} В предыдущем мы пользовались некоторой произвольно
выбранной системой координат и не касались вопроса о законе
преобразования компонент тензора при преобразовании координат. Этот
закон может быть найден из требования (вытекающего из самого определения
понятия натяжения), чтобы слагающие $T_{xn}\,dS$, $T_{yn}\,dS$,
$T_{zn}\,dS$ силы $\mathbf{T}_n\,dS$, действующей на произвольно
расположенную и произвольно ориентированную площадку $dS$,
преобразовывались бы по правилам преобразования
векторов\,\footnote{Однако, например, величина $\mathbf{T}_x$ со
слагающими $T_{xx}$, $T_{yx}$, $T_{zx}$, определяющая силу,
действующую на некоторую площадку, сама не является вектором, ибо
само направление площадки зависит от произвольно выбранного
направления координатной оси $x$.}.

Мы не будем останавливаться здесь на выводе этого закона
преобразования; отметим только, что с помощью его можно убедиться
в том, что как уравнение (33.7), так и условие (33.10) симметрии
тензора сохраняют свой вид при любом преобразовании декартовых
координат.

\section*{\S\,34. Тензор натяжений электрического поля.}

\textbf{1.} Обращаемся к поставленной в начале предыдущего параграфа
задаче сведения пондеромоторных сил электрического поля к натяжениям.
С этой целью удобно разложить общее выражение (32.10) для объёмной
плотности этих сил на два слагаемых:
\begin{equation}
    \mathbf{f} = \mathbf{f}' + \mathbf{f}'',
    \quad
    \mathbf{f}' = \rho\mathbf{E} - \frac{1}{8\pi}E^2\operatorname{grad}\varepsilon,
    \quad
    \mathbf{f}'' = \frac{1}{8\pi}\operatorname{grad}\!\left(
        E^2\frac{\partial\varepsilon}{\partial\tau}\tau
    \right).
    \tag{34.1}
\end{equation}
Наша задача будет, очевидно, разрешена, если мы найдём такой тензор
$\mathbf{t}$, чтобы по подстановке его компонент в правую часть
уравнений (33.7) левые части этих уравнений совпали с определяемыми
формулой (34.1) слагаемыми плотности объёмных сил $\mathbf{f}'$.

\textbf{2.} Выразим в (34.1) $\rho$ через $\operatorname{div}\mathbf{D}$
с помощью (22.2) и рассмотрим слагающую плотности сил $\mathbf{f}'$
по какому-нибудь направлению, например по оси $x$:
$$
    f_x = \frac{1}{4\pi}E_x\operatorname{div}\mathbf{D}
        - \frac{1}{8\pi}E^2\frac{\partial\varepsilon}{\partial x}.
$$

% стр. 162 (левая колонка)

Легко убедиться [ср.\ уравнение $(43^*_a)$, где $\varphi$ соответствует
в нашем случае величине $E_x$], что
$$
    E_x\operatorname{div}\mathbf{D}
    = \frac{\partial}{\partial x}(E_x D_x)
    + \frac{\partial}{\partial y}(E_x D_y)
    + \frac{\partial}{\partial z}(E_x D_z)
    - \mathbf{D}\nabla\cdot E_x.
$$
Далее, воспользовавшись (32.2), имея в виду (7.6), получаем:
$$
    \mathbf{D}\nabla\cdot E_x = \varepsilon\mathbf{E}\nabla\cdot E_x
    = \frac{\varepsilon}{2}\frac{\partial}{\partial x}E^2,
$$
и, следовательно,
$$
    -\frac{1}{4\pi}\mathbf{D}\nabla\cdot E_x
    - \frac{1}{8\pi}E^2\frac{\partial\varepsilon}{\partial x}
    = -\frac{1}{8\pi}\left\{
        \varepsilon\frac{\partial E^2}{\partial x}
      + E^2\frac{\partial\varepsilon}{\partial x}
    \right\}
    = -\frac{1}{8\pi}\frac{\partial(\varepsilon E^2)}{\partial x}.
$$
Таким образом,
$$
    f_x = \frac{1}{4\pi}\left\{
        \frac{\partial}{\partial x}\!\left(E_x D_x - \frac{\varepsilon E^2}{2}\right)
      + \frac{\partial}{\partial y}(E_x D_y)
      + \frac{\partial}{\partial z}(E_x D_z)
    \right\}.
$$
Уравнение это совпадёт по форме с (33.7), если положить:
$$
    T'_{xx} = \frac{\varepsilon}{4\pi}\!\left(E_x^2 - \frac{E^2}{2}\right),
    \quad
    T'_{xy} = \frac{\varepsilon}{4\pi}E_x E_y,
    \quad
    T'_{xz} = \frac{\varepsilon}{4\pi}E_x E_z.
$$
Аналогичным образом могут быть определены и остальные компоненты
тензора $\mathbf{T}'$, совокупность которых может быть представлена
в форме таблицы типа (33.6):
\begin{equation}
    4\pi\mathbf{T}' =
    \begin{pmatrix}
        \varepsilon\!\left(E_x^2 - \dfrac{E^2}{2}\right)
            & \varepsilon E_x E_y & \varepsilon E_x E_z \\
        \varepsilon E_y E_x
            & \varepsilon\!\left(E_y^2 - \dfrac{E^2}{2}\right)
            & \varepsilon E_y E_z \\
        \varepsilon E_z E_x
            & \varepsilon E_z E_y
            & \varepsilon\!\left(E_z^2 - \dfrac{E^2}{2}\right)
    \end{pmatrix}.
    \tag{34.2}
\end{equation}
Что же касается второй слагающей $\mathbf{f}''$ объёмных сил, то, как
непосредственно видно, она совпадает по форме с (33.7), причём отличны
от нуля только те компоненты эквивалентного ей тензора натяжений
$\mathbf{T}''$, которые расположены на главной диагонали таблицы (33.6):
\begin{equation}
    T''_{xx} = T''_{yy} = T''_{zz}
    = \frac{1}{8\pi}E^2\frac{\partial\varepsilon}{\partial\tau}\tau,
    \quad
    T_{lk} = 0 \text{ при } l \neq k.
    \tag{34.3}
\end{equation}
Таким образом, как $\mathbf{T}'$, так и $\mathbf{T}''$ (а стало быть,
и их сумма $\mathbf{T} = \mathbf{T}' + \mathbf{T}''$) являются
симметричными тензорами, т.\,е. удовлетворяют условию (33.10).

Итак, эквивалентность объёмных сил (34.1) системе натяжений (34.2)
и (34.3) нами доказана, и мы можем утверждать, что общая сила $F$,
действующая на произвольный участок среды $V$, определяется состоянием
поля \textit{на границах этого участка}, т.\,е. может быть сведена
к системе сил или натяжений, приложенных к его поверхности $S$.

% стр. 163 (правая колонка)

\textbf{3.} Впервые пондеромоторные силы поля были сведены к натяжениям
Максвеллом, который, однако, не учитывал зависимости диэлектрической
постоянной $\varepsilon$ от плотности диэлектрика (см.\ конец §\,32).
Поэтому \textit{максвеллов тензор натяжений} соответствует лишь части
нашего тензора натяжений $\mathbf{T}$, а именно тензору $\mathbf{T}'$.

Для отличия от «максвелловских» сил $\mathbf{f}'$ и «максвелловских»
натяжений $\mathbf{T}'$ мы будем называть силы $\mathbf{f}''$ и
соответствующий им тензор $\mathbf{T}''$ \textit{стрикционными}.
Заметим, что стрикционные натяжения $\mathbf{T}''$ эквивалентны
всестороннему давлению\,\footnote{В твёрдых диэлектриках наши выражения
для стрикционных сил и натяжений, учитывающие лишь объёмную стрикцию,
должны быть дополнены членами, величина которых определяется
зависимостью $\varepsilon$ от деформаций, не сопровождающихся
изменениями плотности среды.}
\begin{equation}
    p'' = -\frac{1}{8\pi}E^2\frac{\partial\varepsilon}{\partial\tau}\tau,
    \tag{34.4}
\end{equation}
ибо давление есть отрицательное натяжение, направленное по нормали
к произвольной площадке внутри тела и не зависящее от ориентации этой
площадки.

В жидких и твёрдых телах стрикционные силы и натяжения одного порядка
величины с максвелловскими силами и натяжениями
$\left(\text{ибо }\frac{\partial\varepsilon}{\partial\tau}\tau\text{ одного порядка с }\varepsilon\right)$.
Это отмечалось Гельмгольцем, Джинсом и др., однако оставалось неясным,
почему пренебрегающие стрикционными силами и натяжениями авторы не
получали существенно ошибочных результатов. Покажем, в чём здесь дело.

Допустим, что мы интересуемся только равнодействующей $\mathbf{F}$
всех сил электрического поля, приложенных к произвольному телу $A$,
и результирующим моментом $\mathbf{N}$ этих сил. Согласно (33.3)
и (33.9), равнодействующая $\mathbf{F}$ и момент $\mathbf{N}$ этих сил
однозначно определяются значениями тензора $\mathbf{T}$
\textit{с внешней стороны поверхности тела} $S$. Стрикционные натяжения
$\mathbf{T}''$ в вакууме, очевидно, равны нулю (ибо в вакууме
$\frac{\partial\varepsilon}{\partial\tau}\tau$ нужно считать равным
нулю). Поэтому, если тело $A$ окружено вакуумом, то стрикционные силы
$\mathbf{f}''$ и натяжения $\mathbf{T}''$ влияют только на распределение
сил по объёму тела $A$\,\footnote{Этим распределением определяются
возникающие в теле под воздействием электрического поля упругие
натяжения — так называемая \textit{электрострикция}, откуда и название
сил $\mathbf{f}''$ и натяжений $\mathbf{T}''$.}, но не влияют ни на
величину равнодействующей всех сил $\mathbf{F}$, ни на их момент
$\mathbf{N}$.

Если же тело $A$ окружено не вакуумом, а диэлектриком $B$, то
определяемая тензором $\mathbf{T}''$ равнодействующая $\mathbf{F}''$
стрикционных сил $\mathbf{f}''$, вообще говоря, не равна нулю. Однако,
если система тел $A$ и $B$ находится в механическом равновесии, то эта
не учитывавшаяся Максвеллом сила $\mathbf{F}''$ точно компенсируется
также не учитывавшимися Максвеллом дополнительными \textit{механическими}
















% ============================================================
% ПРОДОЛЖЕНИЕ §34 (стр. 164–168)
% ============================================================

% стр. 164 (левая колонка — продолжение §34, п.3)

силами, испытываемыми телом $A$ со стороны тела $B$ благодаря
тому, что само тело $B$ подвергается электрострикции.

Поясним это простейшим примером. Пусть тело $A$ представляет
собой маленький заряженный шарик, погружённый в бесконечный
однородный жидкий диэлектрик $B$; гидростатическое давление
в жидкости $B$ вдали от шарика $A$ задано и равно $p_B^0$. Весом
жидкости $B$ пренебрегаем. Тогда «по Максвеллу» на поверхность
$S$ шарика $A$ будут действовать, во-первых, электрические натяжения
$\mathbf{T}'$ и, во-вторых, гидростатическое давление $p_B^0$.
В действительности же надо ещё учесть, во-первых, электрострикционные
натяжения $\mathbf{T}''$, действующие на поверхность тела $A$
и эквивалентные, согласно (34.4), давлению $p''$; во-вторых, так
как жидкость подвергается электрострикции в электрическом поле
давлению $p''$, то гидростатическое давление $p_B$ жидкости вблизи
шарика $A$ не будет равно гидростатическому давлению $p_B^0$
в бесконечности: условием равновесия жидкости будет постоянство
полного давления в ней
$$
    p'' + p_B = \mathrm{const} = p_B^0.
$$

Таким образом, сумма действующих на поверхность тела электрострикционных
натяжений $\mathbf{T}''$, эквивалентных гидростатическому давлению
$p''$, и фактического гидростатического давления жидкости $p_B$ равна
гидростатическому давлению жидкости $p_B^0$, вычисленному без учёта
электрострикции\,\footnote{Строго говоря, надо учесть, что благодаря
электрострикции изменяется плотность $\tau$, а вместе с ней и
диэлектрическая постоянная $\varepsilon$ жидкости. Благодаря этому
изменяются поле заряженного шарика $A$ в жидкости и действующие на него
максвелловы натяжения $\mathbf{T}'$. Однако, как показывает подсчёт,
эти изменения натяжений совершенно ничтожны.}. Конечно, если тело $A$
не находится в равновесии, то неучёт электрострикции поведёт, вообще
говоря, к ошибке.

Итак, если нас интересует не распределение пондеромоторных сил
по объёму произвольного тела $A$, а лишь равнодействующая $\mathbf{F}$
этих сил и их момент $\mathbf{N}$, то можно ограничиться рассмотрением
одних только максвелловских сил $\mathbf{f}'$ и максвелловских натяжений
$\mathbf{T}'$, отбрасывая стрикционные силы и натяжения $\mathbf{f}''$
и $\mathbf{T}''$, при условии, что тело $A$ окружено либо вакуумом,
либо диэлектрической средой, находящейся в механическом равновесии.

\textbf{4.} Внося в (33.5) значения компонент тензора $\mathbf{T}$,
можно определить слагающую по оси $x$ силы натяжения $\mathbf{T}_n$,
действующей в электрическом поле на единичную площадку, внешняя
нормаль к которой направлена по $\mathbf{n}$:
\begin{align*}
    T_{xn} &= \frac{1}{4\pi}\left\{\varepsilon E_x^2
            - \frac{1}{2}\left(\varepsilon - \frac{\partial\varepsilon}{\partial\tau}\tau\right)E^2
            \right\}\cos(\mathbf{n},x) + \\
           &\quad + \frac{\varepsilon}{4\pi}E_x E_y\cos(\mathbf{n},y)
            + \frac{\varepsilon}{4\pi}E_x E_z\cos(\mathbf{n},z) = \\
           &= \frac{\varepsilon}{4\pi}E_x E_n
            - \frac{1}{8\pi}E^2\!\left(\varepsilon
            - \frac{\partial\varepsilon}{\partial\tau}\tau\right)\cos(\mathbf{n},x).
\end{align*}

% стр. 165 (правая колонка)

Совокупность этого выражения для $T_{xn}$ и аналогичных выражений
для $T_{yn}$ и $T_{zn}$ в векторной форме может быть записана так:
\begin{equation}
    \mathbf{T}_n = \frac{\varepsilon}{4\pi}E_n\mathbf{E}
                 - \frac{1}{8\pi}\mathbf{n}E^2\!\left(
                     \varepsilon - \frac{\partial\varepsilon}{\partial\tau}\tau
                   \right).
    \tag{34.5}
\end{equation}

В частном случае, если рассматриваемый элемент поверхности
перпендикулярен полю $\mathbf{E}$ и его внешняя нормаль $\mathbf{n}$
параллельна $\mathbf{E}$, натяжение, приложенное к нему извне,
определится выражением
\begin{equation}
    \mathbf{T}_n = \frac{\varepsilon + \dfrac{\partial\varepsilon}{\partial\tau}\tau}{8\pi}
                   E^2\mathbf{n}
    \qquad (\mathbf{n} \parallel \mathbf{E});
    \tag{34.6}
\end{equation}

если же нормаль $\mathbf{n}$ перпендикулярна полю, то
$$
    E_n = 0
$$
и
\begin{equation}
    \mathbf{T}_n = -\frac{\varepsilon - \dfrac{\partial\varepsilon}{\partial\tau}\tau}{8\pi}
                   E^2\mathbf{n}
    \qquad (\mathbf{n} \perp \mathbf{E}).
    \tag{34.7}
\end{equation}

Наконец, если площадка $dS$ имеет некоторое промежуточное направление,
то её можно разложить на две взаимно перпендикулярные площадки,
параллельную и перпендикулярную $\mathbf{E}$, причём действующие на
эти площадки силы определяются из (34.6) и (34.7). Стало быть, система
натяжений в электрическом поле сводится к \textit{тяге}
$\dfrac{\varepsilon + \frac{\partial\varepsilon}{\partial\tau}\tau}{8\pi}E^2$
\textit{по направлению поля} $\mathbf{E}$ [уравнение (34.6)] и к
\textit{давлению} (отрицательная тяга)
$\dfrac{\varepsilon - \frac{\partial\varepsilon}{\partial\tau}\tau}{8\pi}E^2$
\textit{по направлению, перпендикулярному} $\mathbf{E}$ [уравнение (34.7)].

С точки зрения механистической теории поля, эта тяга и это давление
суть не что иное, как силы упругости, возникающие в эфире при
деформации его в электрическом поле. Согласно (34.6) и (34.7), силы
эти испытываются всеми теми участками эфира (как в вакууме, так и
в материальных телах), в которых поле (т.\,е. деформация эфира) не
равно нулю\,\footnote{Мы получим правильную картину этих натяжений,
если представим себе, что вдоль силовых линий поля натянуты
материальные упругие нити, подверженные тяге
$\dfrac{\varepsilon + \frac{\partial\varepsilon}{\partial\tau}\tau}{8\pi}E^2$
и оказывающие друг на друга боковое давление
$\dfrac{\varepsilon - \frac{\partial\varepsilon}{\partial\tau}\tau}{8\pi}E^2$.}.

% стр. 166 (левая колонка)

Конечно, с точки зрения современной теории, отрицающей существование
материального эфира (в механистическом смысле этого слова),
пондеромоторные силы электромагнитного поля могут быть приложены лишь
к электрическим зарядам и к материальным телам, несущим эти заряды,
или, точнее, состоящим из электрических зарядов (электронов и атомных
ядер). Однако, по доказанному, результирующая сила, действующая на
тела, находящиеся в произвольном объёме $V$, может быть формально
представлена в виде суммы натяжений, «испытываемых» поверхностью этого
объёма $S$ (могущей, конечно, проходить как в вакууме, так и
в материальных телах). Следовательно, мы можем оперировать с этими
натяжениями, будучи уверенными в правильности окончательных результатов.
Принципиальное же значение понятия натяжений электромагнитного поля
выяснится в §\,105, в котором мы убедимся, что в \textit{переменных}
электромагнитных полях нарушается эквивалентность между пондеромоторными
силами и электромагнитными натяжениями и что избыток суммы натяжений,
испытываемых поверхностью объёма $V$, над пондеромоторными силами,
испытываемыми находящимися в этом объёме телами, определяет собою
изменение количества движения электромагнитного поля в этом объёме.

\textbf{5.} Замена пондеромоторных сил эквивалентной системой натяжений
весьма облегчает, в частности, определение сил, приложенных к
поверхностям разрыва электрического поля, т.\,е. к поверхностям,
к которым подходят поля с различными значениями $\varepsilon$, и к
поверхностям раздела заряженной свободным электричеством поверхности
различной поляризуемости. Правда, при выводе системы натяжений (34.2)
и (34.3) мы базировались на результатах §\,32, в котором предполагались
отсутствующими подобные поверхности разрыва. Однако, определив на
основании (34.2) и (34.3) результирующую силу, приложенную, например,
к слою конечной толщины, и переходя затем к пределу к бесконечно
тонкому слою, мы определим, очевидно, величину силы, действующей на
заряженную поверхность.

Нормаль к поверхности слоя, отграничивающего среду $1$ от среды $2$,
обозначим через $\mathbf{n}$; для определённости предположим, что
$\mathbf{n}$ направлено из среды $1$ в среду $2$. Элемент $dS$ той
поверхности слоя, которая граничит со средой $1$, будет, согласно
(34.5), испытывать со стороны этой среды силу
$$
    \mathbf{T}'_{1n}\,dS = -\frac{\varepsilon_1}{8\pi}
    (2E_{1n}\mathbf{E}_1 - E_1^2\mathbf{n})\,dS
$$
(для простоты тензор $\mathbf{T}$ заменяем максвелловым тензором
$\mathbf{T}'$), а элемент $dS$ поверхности, граничащей со средой $2$,
будет испытывать со стороны этой среды силу
$$
    \mathbf{T}'_{2n}\,dS = +\frac{\varepsilon_2}{8\pi}
    (2E_{2n}\mathbf{E}_2 - E_2^2\mathbf{n})\,dS.
$$
Отличие знаков этих выражений обусловливается тем, что, по условию,
направление \textit{внешней} нормали к слою совпадает в среде $2$
с направлением $\mathbf{n}$,

% стр. 167 (правая колонка)

а в среде $1$ — прямо ему противоположно. Общая сила, действующая
на элемент слоя поверхности $dS$, равна, таким образом:
\begin{align*}
    \mathbf{f}'\,dS &= (\mathbf{T}'_{1n} + \mathbf{T}'_{2n})\,dS = \\
    &= \frac{1}{4\pi}(D_{2n}\mathbf{E}_2 - D_{1n}\mathbf{E}_1)
     - \frac{1}{8\pi}(D_2 E_2 - D_1 E_1)\mathbf{n}\,dS.
\end{align*}
Очевидно, что это соотношение должно остаться справедливым и при
переходе к предельному случаю бесконечно тонкого слоя, т.\,е.
к поверхности. Таким образом, результирующая сила, действующая на
единицу площади произвольной поверхности, равна:
\begin{equation}
    \mathbf{f}' = \frac{1}{4\pi}(D_{2n}\mathbf{E}_2 - D_{1n}\mathbf{E}_1)
               - \frac{1}{8\pi}(D_2 E_2 - D_1 E_1)\mathbf{n},
    \tag{34.8}
\end{equation}
где $\mathbf{E}_1$, $\mathbf{D}_1$ и $\mathbf{E}_2$, $\mathbf{D}_2$ суть,
соответственно, значения векторов $\mathbf{E}$ и $\mathbf{D}$ с внутренней
и внешней (по отношению к $\mathbf{n}$) сторон этой поверхности.

Разумеется, для всякой поверхности, не являющейся поверхностью разрыва,
$\mathbf{E}_1 = \mathbf{E}_2$ и $\mathbf{D}_1 = \mathbf{D}_2$, так что
сила $\mathbf{f}'$ обращается в нуль.

Приложенные к поверхности разрыва силы (34.8) проявляются в том, что
под их воздействием в помещённых в электрическое поле телах возникают
уравновешивающие их \textit{упругие} натяжения; сумма всех натяжений
(электрических и упругих) не может испытывать разрывов и должна быть
одинаковой по обе стороны любой поверхности, в том числе и поверхности
заряженной или поверхности раздела двух сред различной поляризуемости.
Впрочем, в большинстве случаев представляют интерес не натяжения,
возникающие в телах под воздействием электрического поля, а
результирующая сила, действующая на данное тело и определяющаяся не
скачком тензора натяжений $\mathbf{T}$ на поверхности тела $S$, а
значением слагающих тензора $\mathbf{T}$ с \textit{внешней стороны
поверхности тела} $S$.

\textbf{6.} Воспользуемся теперь тензором натяжений $\mathbf{T}$, чтобы
уточнить смысл обычного утверждения, гласящего: пробный заряд $e$,
помещённый в жидкий или газообразный диэлектрик, напряжённость
электрического поля в котором равна $\mathbf{E}_0$, испытывает силу
\begin{equation}
    \mathbf{F} = e\mathbf{E}_0.
    \tag{34.9}
\end{equation}

Пусть пробный заряд $e$ представляет собою небольшое заряжённое
произвольной формы тело $A$ из металла или диэлектрика. Обозначим через
$\mathbf{E}'$ поле этого тела, помещённого в данную точку среды, так
что истинное поле $\mathbf{E} = \mathbf{E}_0 + \mathbf{E}'$. Поле
$\mathbf{E}'$ возбуждается, во-первых, свободным зарядом $e$ пробного
тела $A$ и, во-вторых, распределением связанных или индуцированных
зарядов, которые возникают в нём под воздействием внешнего поля
$\mathbf{E}_0$. Результирующую силу, действующую на пробное тело $A$,
можно вычислить с помощью формулы (34.2); тензор $\mathbf{T}''$, по
доказанному, к этой силе ничего не прибавляет.

Слагающие тензора $\mathbf{T}'$ являются квадратичными функциями
слагающих поля $\mathbf{E}$. При вычислении слагающих $\mathbf{T}'$
делаем два предположения:

1) тело $A$ настолько мало, что на всём протяжении некоторой
охватывающей его поверхности $S$ невозмущённое наличием пробного тела
$\mathbf{E}_0$ и диэлектрическую постоянную среды $\varepsilon$ можно
считать постоянными;

2) заряд $e$ и размеры тела $A$ настолько малы, что возбуждаемое им
в точках этой поверхности $S$ поле $\mathbf{E}'$ гораздо меньше
$\mathbf{E}_0$\,\footnote{Заметим, что размеры поверхности $S$ могут
значительно превышать размеры самого пробного тела и ограничены только
первым условием, ибо на находящиеся внутри $S$ участки однородной среды
силы $\mathbf{f}'$ не действуют.}, так что при вычислении слагающих
$\mathbf{T}'$ членами, квадратичными в $\mathbf{E}'$, можно пренебречь.

% стр. 168 (левая колонка)

Пусть ось $x$ совпадает с направлением поля $\mathbf{E}_0$ вблизи
пробного тела $A$. Тогда при перечисленных условиях путём простых
выкладок получим из (34.2):
$$
    4\pi T'_{xx} = \varepsilon\!\left(\frac{E_0^2}{2} + E_0 E'_x\right),
    \quad
    4\pi T'_{xy} = \varepsilon E_0 E'_y,
    \quad
    4\pi T'_{xz} = \varepsilon E_0 E'_z
$$
и после подстановки в (33.2) и (33.5)
$$
    F_x = F'_x = \frac{\varepsilon}{8\pi}E_0^2\oint\cos(\mathbf{n},x)\,dS
               + \frac{1}{4\pi}E_0\oint \varepsilon E'_n\,dS.
$$
Первый интеграл, как легко убедиться, обращается в нуль\,\footnote{Ср.\
§\,56, в частности формулу (56.3), из которой следует:
$\oint\cos(\mathbf{n},x)\,dS = \oint\mathbf{i}\,dS = \mathbf{i}\oint\,dS = 0$.},
второй же на основании (22.3) равен $4\pi e$, так что $F_x = \varepsilon E_0$.

Далее, при указанных условиях\,\footnote{Действительно, принявши во
внимание, что $\oint\cos(\mathbf{n},y)\,dS = 0$, для $F_y$ легко
получить выражение
$F_y = F'_y = \dfrac{\varepsilon}{4\pi}E_0\oint\{E'_y\cos(x,\mathbf{n})
- E'_x\cos(y,\mathbf{n})\}\,dS$.
Вводя единичные векторы $\mathbf{i}$, $\mathbf{j}$, $\mathbf{k}$
по осям координат, можем написать:
$E'_y\cos(x,\mathbf{n}) - E'_x\cos(y,\mathbf{n})
= (\mathbf{E}'\mathbf{i})(\mathbf{i}\mathbf{n}) - (\mathbf{E}'\mathbf{j})(\mathbf{j}\mathbf{n})
= [\mathbf{i}][\mathbf{n}\mathbf{E}']$.
Наконец, на основании (56*) и (7.6)
$\oint[\mathbf{n}\mathbf{E}']\,dS = \int\operatorname{rot}\mathbf{E}'\,dV = 0$.}
$F_y = F_z = 0$. Таким образом, перечисленные условия действительно
обеспечивают справедливость формулы (34.9).

\medskip
\noindent\textbf{Задача 20.} Заряд $e$ расположен в вакууме на расстоянии
$z_0$ от поверхности однородного диэлектрика, заполняющего полупространство
$z \leqslant 0$. Пользуясь результатами решения примера в конце §\,23,
показать, что сила притяжения между зарядом и диэлектриком равна:
$$
    F = \frac{\varepsilon - 1}{\varepsilon + 1}\left(\frac{e}{2z_0}\right)^2.
$$

\textbf{Примечание.} Конечно, это выражение для силы взаимодействия
остаётся справедливым и в случае диэлектрика конечных размеров, но
только если его поперечные размеры (точнее, расстояние от заряда тех
участков, где поверхность диэлектрика перестаёт быть плоской или где
начинает меняться диэлектрическая постоянная диэлектрика) и толщина
достаточно велики по сравнению с $z_0$.

% ============================================================
% §83. ОПРЕДЕЛЕНИЕ ПОНДЕРОМОТОРНЫХ СИЛ МАГНИТНОГО ПОЛЯ
% ИЗ ВЫРАЖЕНИЯ ЭНЕРГИИ (стр. 386–387)
% ============================================================

\section*{\S\,83. Определение пондеромоторных сил магнитного поля
из выражения энергии\,\footnote{При первом чтении параграф этот может
быть опущен.}.}

% стр. 386 (левая колонка — продолжение §82)

где $\overline{R^3}$ есть средний квадрат расстояния электронов от
ядра атома и, стало быть,
$$
    w_{\text{кин}} = \frac{\overline{N}Ze^2\overline{R^3}}{12mc^2}H^2.
$$

Внося сюда выражение (69.3) для вектора намагничения диамагнетиков,
получаем:
$$
    w_{\text{кин}} = -\frac{1}{2}\mathbf{I}\mathbf{H}.
$$

Вспомним, наконец, что при рассмотрении намагничения диа- и
парамагнетиков в §§\,70 и 71 мы считали, что на атомы магнетика
действует поле $\mathbf{H}$, тогда как правильнее было бы считать,
что на них действует поле $\overline{H}_{\text{микро}} = \mathbf{B}$.
Сделав соответствующую замену в последнем уравнении, получаем
окончательно:
\begin{equation}
    w_{\text{кин}} = -\frac{1}{2}\mathbf{I}\mathbf{B}.
    \tag{82.3}
\end{equation}

Сумма этой плотности кинетической энергии прецессии электронов и
приближённого выражения (82.1) для плотности энергии магнитного поля,
в собственном смысле этого слова\,\footnote{В этом выражении не
учитывается разница между $(\overline{H}_{\text{микро}})^2$ и
$\overline{H^2_{\text{микро}}}$.}, равна [см.\ уравнение (62.5)]:
$$
    w = \frac{1}{8\pi}B^2 - \frac{1}{2}\mathbf{I}\mathbf{B}
      = \frac{1}{8\pi}B(B - 4\pi I) = \frac{1}{8\pi}\mathbf{B}\mathbf{H},
$$
т.\,е., как и требовалось доказать, равна макроскопической плотности
магнитной энергии (82.2).

\bigskip
\noindent\textbf{\S\,83. Определение пондеромоторных сил магнитного поля
из выражения энергии.}

\medskip
\textbf{1.} В §\,66 мы определили пондеромоторные силы, испытываемые
магнетиками в магнитном поле, исходя из определённых представлений
о молекулярном строении магнетиков. Теперь мы дадим более общий вывод
выражения пондеромоторных сил магнитного поля, исходя из энергии
магнитного поля (81.3) и рассматривая изменение этой энергии $\delta W$
при бесконечно малом виртуальном перемещении $\mathbf{q}$ находящихся
в поле тел. Если при этих виртуальных перемещениях оставлять силы токов
проводимости постоянной, то, согласно уравнению (79.9), изменение
магнитной энергии при этих перемещениях $\delta W$ будет равно работе
$A$ механических сил:
$$
    (\delta W)_J = A.
$$

% стр. 387 (правая колонка)

С другой стороны, работа эта, очевидно, равна:
$$
    A = \int \mathbf{q}\mathbf{f}\,dV,
$$
где $\mathbf{f}$ есть плотность пондеромоторных сил. Следовательно,
\begin{equation}
    (\delta W)_J = \int \mathbf{q}\mathbf{f}\,dV.
    \tag{83.1}
\end{equation}

Вычислив $(\delta W)_J$ из этого соотношения, можно, очевидно,
определить и искомое значение плотности сил $\mathbf{f}$.

Энергия магнитного поля может быть выражена одним из следующих
уравнений (81.3) и (81.1):
$$
    W_1 = \frac{1}{8\pi}\int \mathbf{H}\mathbf{B}\,dV
        = \frac{1}{8\pi}\int \frac{1}{\mu}B^2\,dV,
$$
$$
    W_2 = W_1 = \frac{1}{2c}\int \mathbf{A}\mathbf{j}\,dV.
$$

Эти выражения энергии, как неоднократно указывалось, справедливы лишь
\textit{в отсутствии ферромагнетиков}; следовательно, все наши выводы
будут применимы лишь к диа- и парамагнитным средам.

Все дальнейшие выкладки весьма аналогичны вычислениям, проведённым
нами при определении пондеромоторных сил электрического поля.
В частности, мы предположим, что поверхностей разрыва в поле нет и
что объём интегрирования охватывает полное поле, так что все
поверхностные интегралы, с которыми мы будем встречаться в дальнейшем,
обратятся в нуль.

\textbf{2.} Изменение энергии поля при произвольном бесконечно малом
перемещении $\mathbf{q}$ находящихся в нём тел равно:
$$
    \delta W_1 = \frac{1}{8\pi}\int B^2\delta\!\left(\frac{1}{\mu}\right)dV
               + \frac{1}{4\pi}\int \frac{1}{\mu}\mathbf{B}\,\delta\mathbf{B}\,dV,
$$
или
$$
    \delta W_2 = \delta W_1 = \frac{1}{2c}\int \mathbf{A}\,\delta\mathbf{j}\,dV
               + \frac{1}{2c}\int \mathbf{j}\,\delta\mathbf{A}\,dV,
$$
где $\delta\!\left(\dfrac{1}{\mu}\right)$, $\delta\mathbf{j}$,
$\delta\mathbf{A}$ суть изменения соответствующих величин, обусловленные
перемещением $\mathbf{q}$.

Так как $\mathbf{B} = \operatorname{rot}\mathbf{A}$, то
$\delta\mathbf{B} = \delta\operatorname{rot}\mathbf{A}$. Порядок
выполнения операций дифференцирования и варьирования может быть изменён
без нарушения результата, так что
$$
    \delta\operatorname{rot}\mathbf{A} = \operatorname{rot}(\delta\mathbf{A}).
$$
Следовательно,
$$
    \frac{1}{\mu}\mathbf{B}\cdot\delta\mathbf{B}
    = \mathbf{H}\operatorname{rot}(\delta\mathbf{A}).
$$



% ============================================================
% §83 ПРОДОЛЖЕНИЕ (стр. 388–390)
% ============================================================

% стр. 388 (левая колонка)

Далее, на основании (44*) и (62.7)
$$
    \mathbf{H}\operatorname{rot}(\delta\mathbf{A})
    = \operatorname{div}[\delta\mathbf{A}\cdot\mathbf{H}]
    + \mathbf{H}\cdot\delta\mathbf{A}
    = \operatorname{div}[\delta\mathbf{A}\cdot\mathbf{H}]
    + \frac{4\pi}{c}\mathbf{j}\cdot\delta\mathbf{A}.
$$

Внося это в выражение для $\delta W_1$ и воспользовавшись теоремой
Гаусса (17*), получаем:
$$
    \delta W_1 = \frac{1}{8\pi}\int B^2\delta\!\left(\frac{1}{\mu}\right)dV
               + \frac{1}{4\pi}\oint[\delta\mathbf{A}\cdot\mathbf{H}]_n\,dS
               + \frac{1}{c}\int\mathbf{j}\cdot\delta\mathbf{A}\,dV.
$$

Согласно условиям, формулированным в начале параграфа, поверхностный
интеграл в этом выражении обращается в нуль. Наконец, ввиду того, что
$\delta W_1 = \delta W_2$, приращение энергии $\delta W$ можно, очевидно,
представить также в следующей форме:
$$
    \delta W = 2\delta W_2 - \delta W_1
             = \frac{1}{c}\int\mathbf{A}\,\delta\mathbf{j}\,dV
             - \frac{1}{8\pi}\int B^2\delta\!\left(\frac{1}{\mu}\right)dV.
$$

Таким образом, вычисление $\delta W$ сведено нами к определению изменения
плотности токов проводимости $\mathbf{j}$ и магнитной проницаемости среды
$\mu$ при виртуальном перемещении $\mathbf{q}$ находящихся в поле тел.

\textbf{3.} Локальное изменение величины $1/\mu$ выразится, очевидно,
формулой, аналогичной выражению для $\delta\varepsilon$ [см.\ §\,32, стр.\ 154,
формулу (32.9')]:
$$
    \delta\frac{1}{\mu} = -\frac{\partial}{\partial\tau}\!\left(\frac{1}{\mu}\right)\tau\,
    \operatorname{div}\mathbf{q} - \mathbf{q}\operatorname{grad}\frac{1}{\mu},
$$
где $\tau$ означает плотность среды (массу единицы объёма)\,\footnote{Как
и в случае диэлектриков, мы для простоты пренебрегаем зависимостью
восприимчивости $\mu$ твёрдых магнетиков от деформаций, не связанных
с изменением плотности среды.}.

Что касается виртуального изменения плотности тока $\delta\mathbf{j}$,
то оно может быть определено из того отмеченного в начале параграфа
условия, что при виртуальном перемещении $\mathbf{q}$ сила тока через
произвольную поверхность должна оставаться постоянной (если только эта
поверхность перемещается \textit{вместе} со средой).

Пусть через произвольную поверхность $S$ до смещения $\mathbf{q}$
протекает ток $J$:
$$
    J = \int_S j_n\,dS = \oint_S \mathbf{j}\,dS.
$$

При смещении $\mathbf{q}$ может, во-первых, измениться на $\delta\mathbf{j}$
плотность тока в различных точках поверхности $S$ и, во-вторых, может
деформироваться контур $L$ этой поверхности, благодаря чему увеличится
или уменьшится её общая площадь. Если $ds$ означает элемент контура $L$,
то при смещении $\mathbf{q}$ этот элемент опишет площадку $\delta S =
[\mathbf{q}\,ds]$ (см.\ §\,50, в частности рис. 50). Приращение охватываемой
контуром площади $S$ сводится к сумме площадок $\delta S$,

% стр. 389 (правая колонка)

описанных всеми его элементами $ds$. Следовательно, долженствующее
обращаться в нуль полное изменение тока $J$ через произвольную
испытывающую деформацию поверхность $S$ равно:
$$
    \delta J = \int_S \delta j_n\,dS + \oint_L \mathbf{j}[\mathbf{q}\,ds] = 0.
$$

Преобразуя последний интеграл по теореме Стокса [уравнение (27*)],
получаем:
$$
    \oint_L \mathbf{j}[\mathbf{q}\,ds]
    = \oint_L [\mathbf{j}\mathbf{q}]\,ds
    = \int_S \operatorname{rot}_n[\mathbf{j}\mathbf{q}]\,dS.
$$

Внося это в предшествующее равенство ввиду произвольности поверхности
$S$, получаем окончательно:
\begin{equation}
    \delta\mathbf{j} = -\operatorname{rot}[\mathbf{j}\mathbf{q}].
    \tag{83.2}
\end{equation}

\textbf{4.} Внося приведённые значения $\delta\dfrac{1}{\mu}$ и
$\delta\mathbf{j}$ в выражение для $\delta W$, получаем:
$$
    \delta W = -\frac{1}{c}\int\mathbf{A}\operatorname{rot}[\mathbf{j}\mathbf{q}]\,dV
             + \frac{1}{8\pi}\int B^2\!\left\{
                 \frac{\partial}{\partial\tau}\frac{1}{\mu}\tau\,\operatorname{div}\mathbf{q}
                 + \mathbf{q}\operatorname{grad}\frac{1}{\mu}
               \right\}dV.
$$

На основании уравнения (44*)
$$
    \mathbf{A}\operatorname{rot}[\mathbf{j}\mathbf{q}]
    = \operatorname{div}[[\mathbf{j}\mathbf{q}]\mathbf{A}]
    + [\mathbf{j}\mathbf{q}]\operatorname{rot}\mathbf{A},
$$
причём последний член может быть преобразован следующим образом:
$$
    [\mathbf{j}\mathbf{q}]\operatorname{rot}\mathbf{A}
    = [\mathbf{j}\mathbf{q}]\mathbf{B}
    = -\mathbf{q}[\mathbf{j}\mathbf{B}].
$$

Далее,
$$
    B^2\frac{\partial}{\partial\tau}\frac{1}{\mu}\tau\,\operatorname{div}\mathbf{q}
    = \operatorname{div}\!\left(B^2\frac{\partial}{\partial\tau}\frac{1}{\mu}\tau\mathbf{q}\right)
    - \mathbf{q}\operatorname{grad}\!\left(B^2\frac{\partial}{\partial\tau}\frac{1}{\mu}\tau\right).
$$

Внося эти значения в выражение для $\delta W$ и принимая во внимание,
что объёмный интеграл дивергенции может быть преобразован в интеграл
по поверхности и поэтому, согласно нашему условию, обращается в нуль,
получаем:
$$
    \delta W = \int\mathbf{q}\left\{
        \frac{1}{c}[\mathbf{j}\mathbf{B}]
        - \frac{1}{8\pi}\operatorname{grad}\!\left(B^2\frac{\partial}{\partial\tau}\frac{1}{\mu}\tau\right)
        + \frac{1}{8\pi}B^2\operatorname{grad}\frac{1}{\mu}
    \right\}dV.
$$

Сравнивая это выражение с выражением (83.1) и принявши во внимание,
что $\mathbf{q}$ является произвольной функцией точки, получаем следующее
выражение для плотности пондеромоторных сил:
\begin{equation}
    \mathbf{f} = \frac{1}{c}[\mathbf{j}\mathbf{B}]
               - \frac{1}{8\pi}\operatorname{grad}\!\left(
                   B^2\frac{\partial}{\partial\tau}\frac{1}{\mu}\tau
               \right)
               + \frac{1}{8\pi}B^2\operatorname{grad}\frac{1}{\mu}.
    \tag{83.3}
\end{equation}

% стр. 390 (левая колонка)

Таким образом, пондеромоторные силы слагаются из сил плотности
$\mathbf{f}^{(1)} = \dfrac{1}{c}[\mathbf{j}\mathbf{B}]$, действующих
на обтекаемые током проводники и совпадающих с прежней формулой (65.1),
и из сил плотности
\begin{equation}
    \mathbf{f}^{(2)} = \frac{1}{8\pi}B^2\operatorname{grad}\frac{1}{\mu}
                     - \frac{1}{8\pi}\operatorname{grad}\!\left(
                         B^2\frac{\partial}{\partial\tau}\frac{1}{\mu}\tau
                     \right),
    \tag{83.4}
\end{equation}
действующих на находящиеся в поле магнетики.

\textbf{5.} Нетрудно показать, что плотность этих сил $\mathbf{f}^{(2)}$
совпадает с ранее выведенным выражением (66.6). Последнее выражение, как
отмечалось в §\,66, справедливо лишь для слабо намагничивающихся
магнетиков (ср.\ §\,32). В этом случае, согласно (69.5) и (70.4),
величина $\varkappa/\mu$ пропорциональна числу молекул в единице объёма
среды, т.\,е. пропорциональна плотности среды $\tau$. Стало быть,
$$
    \frac{\varkappa}{\mu} = \frac{\mu - 1}{4\pi\mu} = c\tau,
$$
где $c$ есть некоторая постоянная, от плотности среды не зависящая.
Следовательно,
$$
    \frac{1}{\mu} = 1 - 4\pi c\tau
$$
и
$$
    \tau\frac{\partial}{\partial\tau}\frac{1}{\mu} = -4\pi c\tau = \frac{1}{\mu} - 1.
$$

Внося это в уравнение (83.4), получаем:
\begin{align*}
    \mathbf{f}^{(2)} &= -\frac{1}{8\pi}\operatorname{grad}\!\left\{
        B^2\!\left(\frac{1}{\mu} - 1\right)
    \right\} + \frac{1}{8\pi}B^2\operatorname{grad}\frac{1}{\mu} = \\
    &= \frac{1}{8\pi}\!\left(1 - \frac{1}{\mu}\right)\operatorname{grad}B^2
     = \frac{\mu - 1}{8\pi\mu}\nabla B^2,
\end{align*}
что действительно полностью совпадает с уравнением (66.6).

\textbf{6.} Если в формуле (83.4) заменить $\mathbf{B}$ через $\mu\mathbf{H}$
и выполнить простые преобразования, то она принимает вид
\begin{equation}
    \mathbf{f}^{(2)} = \frac{1}{8\pi}\operatorname{grad}\!\left(
        H^2\tau\frac{\partial\mu}{\partial\tau}
    \right) - \frac{1}{8\pi}H^2\operatorname{grad}\mu,
    \tag{83.5}
\end{equation}
совершенно аналогичный выражению (32.12) пондеромоторных сил,
испытываемых диэлектриками в электрическом поле: формула (83.5)
получается из (32.12) заменой $\mathbf{E}$ на $\mathbf{H}$ и $\varepsilon$
на $\mu$.

В этом обстоятельстве проявляется отмеченный в конце §\,73 факт,
что благодаря эквивалентности элементарных токов и магнитных диполей
вся \textit{макроскопическая} теория магнетиков может быть одинаково
успешно интерпретирована как на основе современной теории, так и с
точки зрения старых теорий магнетизма, исходивших из предположения
о существовании в магнетиках реальных

% стр. 391 (правая колонка)

магнитных зарядов (диполей), т.\,е. из предположения о полном
соответствии между электрическими свойствами диэлектриков и магнитными
свойствами магнетиков. Именно, с точки зрения этих теорий существует
соответствие между векторами $\mathbf{E}$ и $\mathbf{H}$ и между
величинами $\varepsilon$ и $\mu$, тогда как в действительности вектору
$\mathbf{E}$ соответствует не вектор $\mathbf{H}$, а вектор
$\mathbf{B} = \overline{H}_{\text{микро}}$, а диэлектрической постоянной
$\varepsilon$ соответствует не $\mu$, а $1/\mu$.

Отметим, что Максвелл и ряд других авторов (например, Абрагам, Кон
и т.\,д.) не принимали во внимание зависимости магнитной восприимчивости
от плотности среды, благодаря чему выражение пондеромоторных сил
в магнетиках, которым они пользовались, отличалось от (83.5) отсутствием
первого члена, т.\,е. стрикционных сил:
\begin{equation}
    \mathbf{f}'' = \frac{1}{8\pi}\operatorname{grad}\!\left(
        H^2\frac{\partial\mu}{\partial\tau}\tau
    \right).
    \tag{83.6}
\end{equation}

Впрочем, совершенно так же, как и в случае диэлектриков (см.\ §§\,32
и 34), можно показать, что неучитывавшаяся Максвеллом слагающая
$\mathbf{f}''$ плотности сил $\mathbf{f}^{(2)}$ сказывается лишь на
\textit{распределении} пондеромоторных сил по объёму магнетика, тогда
как равнодействующая $\mathbf{F}$ и результирующий момент $\mathbf{N}$
стрикционных сил $\mathbf{f}''$, приложенных ко всем элементам объёма
какого-либо тела, либо равны нулю (если это тело находится в вакууме),
либо уравновешиваются гидростатическим давлением, возникающим под
воздействием магнитного поля в окружающей тело жидкой или газообразной
среде (см.\ также §\,84).

% ============================================================
% §84. ТЕНЗОР НАТЯЖЕНИЯ МАГНИТНОГО ПОЛЯ (стр. 391–394)
% ============================================================

\section*{\S\,84. Тензор натяжения магнитного поля.}

% стр. 391 (правая колонка, продолжение)

\textbf{1.} Пондеромоторные силы магнитного поля совершенно так же,
как и пондеромоторные силы поля электрического, могут быть сведены
к эквивалентной этим силам системе натяжений, характеризуемой некоторым
тензором натяжений $\mathbf{T}$.

Соответствующие вычисления совершенно аналогичны вычислениям §\,34.
Мы представим общее выражение пондеромоторных сил (83.3) как сумму
двух слагаемых:
\begin{equation}
    \mathbf{f} = \mathbf{f}' + \mathbf{f}'', \quad
    \mathbf{f}' = \frac{1}{c}[\mathbf{j}\mathbf{B}]
               - \frac{1}{8\pi}H^2\operatorname{grad}\mu, \quad
    \mathbf{f}'' = \frac{1}{8\pi}\operatorname{grad}\!\left(
        H^2\tau\frac{\partial\mu}{\partial\tau}
    \right)
    \tag{84.1}
\end{equation}
[ср.\ уравнение (83.5)]. Тензор натяжений $\mathbf{T}$ мы также
представим как сумму двух тензоров $\mathbf{T}'$ и $\mathbf{T}''$:
\begin{equation}
    \mathbf{T} = \mathbf{T}' + \mathbf{T}''
    \tag{84.2}
\end{equation}

% стр. 392 (левая колонка)

так, чтобы $\mathbf{T}'$ соответствовало $\mathbf{f}'$, а $\mathbf{T}''$
соответствовало $\mathbf{f}''$, т.\,е. чтобы были удовлетворены уравнения
\begin{equation}
    f'_x = \frac{\partial T'_{xx}}{\partial x}
          + \frac{\partial T'_{xy}}{\partial y}
          + \frac{\partial T'_{xz}}{\partial z}, \quad
    f''_x = \frac{\partial T''_{xx}}{\partial x}
           + \frac{\partial T''_{xy}}{\partial y}
           + \frac{\partial T''_{xz}}{\partial z}
    \tag{84.3}
\end{equation}
и аналогично для $f'_y$, $f''_y$ и т.\,д. [ср.\ уравнение (33.7)].

\textbf{2.} Рассмотрим сначала силы $\mathbf{f}'$ и тензор $\mathbf{T}'$.
Выражая плотность токов $\mathbf{j}$ через $\operatorname{rot}\mathbf{H}$
[уравнение (62.7)], мы можем написать:
$$
    \frac{1}{c}[\mathbf{j}\mathbf{B}]
    = \frac{1}{4\pi}[\operatorname{rot}\mathbf{H}\cdot\mathbf{B}]
    = \frac{\mu}{4\pi}[\operatorname{rot}\mathbf{H}\cdot\mathbf{H}].
$$

Далее, на основании (47*)
$$
    [\operatorname{rot}\mathbf{H}\cdot\mathbf{H}]
    = \mathbf{H}\nabla\cdot\mathbf{H} - \frac{1}{2}\operatorname{grad}H^2
$$
и, стало быть,
\begin{align*}
    \mathbf{f}' &= \frac{\mu}{4\pi}\mathbf{H}\nabla\cdot\mathbf{H}
                 - \frac{\mu}{8\pi}\operatorname{grad}H^2
                 - \frac{1}{8\pi}H^2\operatorname{grad}\mu = \\
                &= \frac{1}{4\pi}\mathbf{B}\nabla\cdot\mathbf{H}
                 - \frac{1}{8\pi}\operatorname{grad}(\mu H^2).
\end{align*}

Рассмотрим слагающую плотности сил $\mathbf{f}'$ по оси $x$:
$$
    f'_x = \frac{1}{4\pi}\mathbf{B}\cdot\nabla H_x
         - \frac{1}{8\pi}\frac{\partial}{\partial x}(\mu H^2).
$$

Легко убедиться [см.\ (43*), где $\varphi$ соответствует в нашем случае
$H_x$], что
$$
    \mathbf{B}\nabla\cdot H_x = \mathbf{B}\cdot\nabla H_x
    = \operatorname{div}(BH_x) - H_x\operatorname{div}\mathbf{B},
$$
или ввиду того, что $\operatorname{div}\mathbf{B} = 0$,
$$
    \mathbf{B}\cdot\nabla H_x
    = \operatorname{div}(\mathbf{B}H_x)
    = \frac{\partial}{\partial x}(B_x H_x)
    + \frac{\partial}{\partial y}(B_y H_x)
    + \frac{\partial}{\partial z}(B_z H_x).
$$

Внося это в предшествующее выражение для $f'_x$, убеждаемся, что
оно совпадает по форме с уравнением (84.3), если положить:
$$
    T'_{xx} = \frac{\mu}{4\pi}\!\left(\mu H_x^2 - \frac{H^2}{2}\right),
    \quad
    T'_{xy} = \frac{\mu}{4\pi}H_x H_y,
    \quad
    T'_{xz} = \frac{\mu}{4\pi}H_x H_z.
$$

Аналогично определяются и остальные компоненты тензора $\mathbf{T}'$,
совокупность которых может быть представлена в форме таблицы типа (33.6):
\begin{equation}
    4\pi\mathbf{T}' = \begin{pmatrix}
        \mu\!\left(H_x^2 - \dfrac{H^2}{2}\right) & \mu H_x H_y & \mu H_x H_z \\
        \mu H_y H_x & \mu\!\left(H_y^2 - \dfrac{H^2}{2}\right) & \mu H_y H_z \\
        \mu H_z H_x & \mu H_z H_y & \mu\!\left(H_z^2 - \dfrac{H^2}{2}\right)
    \end{pmatrix}
    \tag{84.4}
\end{equation}

% стр. 393 (правая колонка)

Таким образом, тензор натяжений магнитного поля $\mathbf{T}'$ может быть
получен из соответствующего тензора $\mathbf{T}'$ электрического поля
[формула (34.2)] простой заменой $\mathbf{E}$ на $\mathbf{H}$ и
$\varepsilon$ на $\mu$\,\footnote{Тот же результат получится и при замене
в (34.2) $\mathbf{E}$ на $\mathbf{B}$ и $\varepsilon$ на $1/\mu$,
например $\varepsilon E_x E_y \to \dfrac{1}{\mu}B_x B_y = \mu H_x H_y$.
С точки зрения современных представлений о поле в магнетиках, именно
последняя замена имеет непосредственный физический смысл (ввиду
соотношений $\mathbf{E} = \overline{E}_{\text{микро}}$,
$\mathbf{B} = \overline{H}_{\text{микро}}$,
$\mathbf{D} = \varepsilon\overline{E}_{\text{микро}}$,
$\mathbf{H} = \dfrac{1}{\mu}\overline{B}_{\text{микро}}$).}. То же самое
относится и к тензору натяжений $\mathbf{T}''$, эквивалентному силам
$\mathbf{f}''$, ибо выражение плотности этих сил в магнитном поле (84.1)
получается из выражения плотности сил $\mathbf{f}''$ в электрическом поле
(34.1) заменой $\mathbf{E}$ на $\mathbf{H}$ и $\varepsilon$ на $\mu$.
Поэтому по аналогии (34.3)
\begin{equation}
    T''_{xx} = T''_{yy} = T''_{zz}
             = \frac{1}{8\pi}H^2\frac{\partial\mu}{\partial\tau}\tau,
    \tag{84.5}
\end{equation}
тогда как недиагональные компоненты тензора $\mathbf{T}''$ ($T''_{xy}$,
$T''_{xz}$ и т.\,д.) равны нулю.

\textbf{3.} Итак, эквивалентность объёмных сил (84.1) системе натяжений
(84.4) и (84.5) нами доказана (заметим, что оба тензора $\mathbf{T}'$ и
$\mathbf{T}''$ симметричны).

Ввиду отмеченной аналогии натяжений магнитного поля с натяжениями
электрического поля все результаты, полученные нами в §\,34, непосредственно
применимы и к натяжениям магнитного поля. Так, например, система натяжений
в магнитном поле сводится к \textit{тяге}
$\dfrac{\mu + \tau\frac{\partial\mu}{\partial\tau}}{8\pi}H^2$
\textit{по направлению поля} $\mathbf{H}$ и к \textit{давлению}
$\dfrac{\mu - \tau\frac{\partial\mu}{\partial\tau}}{8\pi}H^2$
\textit{по направлению, перпендикулярному к} $\mathbf{H}$
[см.\ уравнения (34.6) и (34.7)].

Далее, так же как и в §\,34, можно показать, что равнодействующая
$\mathbf{F}$ и результирующий момент $\mathbf{N}$ сил, испытываемых
в магнитном поле каким-либо телом, полностью определяются
\textit{максвелловым тензором натяжений} $\mathbf{T}'$ и не зависят от
неучитывавшегося Максвеллом тензора $\mathbf{T}''$, так что стрикционные
натяжения $\mathbf{T}''$ влияют лишь на \textit{распределение}
пондеромоторных сил по объёму тела, а также обусловливают возникновение
в окружающей тело жидкой или газообразной среде уравновешивающего их
гидростатического давления [см.\ уравнение (34.4)]
\begin{equation}
    p = -\frac{1}{8\pi}H^2\frac{\partial\mu}{\partial\tau}\tau.
    \tag{84.6}
\end{equation}

% стр. 394 (левая колонка)

\textbf{4.} Хотя полученные в §§\,81–83 выражения для плотности энергии
и плотности объёмных сил магнитного поля, как неоднократно отмечалось,
применимы лишь в неферромагнитных средах, однако сведение объёмных сил
к натяжениям позволяет, как было показано в §\,75, определить также
равнодействующую $\mathbf{F}$ и результирующий момент $\mathbf{N}$ сил,
действующих на ферромагнетики (но не распределение сил по объёму
ферромагнетика).

% ============================================================
% §85. ВИХРИ ЭЛЕКТРИЧЕСКОГО ПОЛЯ (стр. 394–395)
% ============================================================

\section*{\S\,85. Вихри электрического поля.}

% стр. 394 (левая колонка, продолжение)

\textbf{1.} В §\,76 мы вывели законы индукции токов в \textit{движущихся}
проводниках, основываясь на том, что, согласно §\,45, на электрические
заряды действует лоренцова сила (45.4):
$$
    \mathbf{F} = e\!\left\{\mathbf{E} + \frac{1}{c}[\mathbf{v}\mathbf{H}]\right\},
$$
второй член которой пропорционален скорости заряда и напряжённости
магнитного поля $\mathbf{H}$. Затем, основываясь на принципе
относительности движения, мы показали, что индукция токов должна иметь
место и в \textit{неподвижных} проводниках при изменениях магнитного
поля, причём
\begin{equation}
    \mathcal{E}^{\text{инд}} = -\frac{1}{c}\frac{d\Psi}{dt}
    = -\frac{1}{c}\frac{d}{dt}\int B_n\,dS,
    \tag{85.1}
\end{equation}
где интегрирование может быть распространено по любой поверхности,
опирающейся на контур проводника. Для случая \textit{неподвижных}\,\footnote{Неподвижных
относительно инерциальной системы, в которой производятся измерения
поля (см.\ §\,77).} проводников поверхность эта тоже может быть выбрана
неподвижной, причём в этом случае дифференцирование по времени может
быть выполнено под знаком интеграла:
$$
    \mathcal{E}^{\text{инд}} = -\frac{1}{c}\int\frac{\partial B_n}{\partial t}\,dS.
$$

Знак полной производной по времени заменён нами знаком частной
производной (круглое $\partial$) для того, чтобы отметить, что
$\dfrac{\partial B_n}{\partial t}$ есть скорость изменения во времени
величины $B_n$ в \textit{фиксированной} точке пространства.

Итак, мы приходим к заключению, что изменение магнитного поля должно
вызывать в неподвижных проводниках появление сил, действующих на
электрические заряды, причём циркуляция этих сил по контуру проводника,
обозначаемая нами через $\mathcal{E}^{\text{инд}}$, определяется
формулой (85.1).

\textbf{2.} В §\,2 напряжённость электрического поля $\mathbf{E}$ была
определена нами как сила, действующая на единичный положительный пробный

% стр. 395 (правая колонка)

заряд. Однако в §\,45 мы убедились, что и в отсутствии электрического
поля \textit{движущийся} заряд может испытывать силу
$$
    \mathbf{F} = \frac{e}{c}[\mathbf{v}\mathbf{H}].
$$

Это обстоятельство ведёт к необходимости уточнить определение
напряжённости электрического поля $\mathbf{E}$ в том смысле, что
$\mathbf{E}$ равно силе, действующей на \textit{неподвижный} единичный
положительный заряд. Действительно, из уравнения (45.4) следует при
$\mathbf{v} = 0$:
$$
    \mathbf{E} = \frac{1}{e}\mathbf{F}
$$
(предполагаем, что сторонние электродвижущие силы химического и
термического происхождения отсутствуют).

Исходя из этого определения электрического поля, мы на основании
относящейся к \textit{неподвижным} проводникам формулы (85.1) должны
заключить, что при изменениях магнитного поля в этих проводниках
возбуждается поле электрическое, циркуляция напряжённости которого по
контуру проводника $L$ равна:
\begin{equation}
    \oint_L E_s\,ds = \mathcal{E}^{\text{инд}}
    = -\frac{1}{c}\int\frac{\partial B_n}{\partial t}\,dS
    = -\frac{1}{c}\frac{\partial\Psi}{\partial t}.
    \tag{85.2}
\end{equation}

Таким образом, явления индукции токов в проводниках, \textit{движущихся}
в постоянном магнитном поле, истолковываются нами как результат
воздействия \textit{магнитного} поля (лоренцова сила), тогда как
индукция в \textit{неподвижных} проводниках при изменениях магнитного
поля истолковывается совсем иным образом — как результат воздействия
\textit{электрического} поля, возбуждаемого изменениями поля магнитного.
Между тем, как мы убедились в §\,76, никакой объективной разницы между
этими двумя видами индукции нет, ибо понятие движения относительно.
Первая из посвящённых теории относительности работ Эйнштейна начинается
с указанием на необходимость устранения этой принципиальной разницы
в истолковании двух явлений, объективно неотличимых друг от друга.
Теория относительности эту задачу разрешила, о чём вкратце будет
рассказано в §\,117.

\textbf{3.} Так как, согласно уравнению (7.3), циркуляция электрического
вектора поля стационарных зарядов равна нулю, то формула (85.2) остаётся
справедливой и в том случае, если мы условимся во всём дальнейшем
понимать под $\mathbf{E}$ общую напряжённость электрического поля вне
зависимости от того, возбуждается ли это поле (полностью или частично)
\textit{стационарными электрическими зарядами} (кулоново поле) \textit{или
же изменениями поля магнитного}. При выводе уравнения (85.2) предполагалось,
что контур интегрирования $L$ совпадает с контуром линейного проводника.
Естественно, однако, предположить, что если изменения магнитного

% ============================================================
% §104. ЭЛЕКТРОМАГНИТНЫЙ МОМЕНТ КОЛИЧЕСТВА ДВИЖЕНИЯ.
% ЧАСТНЫЙ СЛУЧАЙ СТАТИЧЕСКОГО ПОЛЯ (стр. 502–506)
% ============================================================

\section*{\S\,104. Электромагнитный момент количества движения.
Частный случай статического поля.}

% стр. 502 (левая колонка)

\textbf{1.} Как понятие электромагнитного количества движения $\mathbf{g}$,
так и понятие потока электромагнитной энергии (вектора Пойнтинга
$\mathbf{S}$) было сформулировано на основе изучения переменных (в
частности, волновых) полей. Посмотрим теперь на частном примере, к каким
выводам приводит применение этих понятий к полям статическим.

Рассмотрим цилиндрический конденсатор, помещённый в однородное
магнитное поле $\mathbf{H}$, параллельное его оси (рис. 86). В
пространстве между обкладками конденсатора, помимо магнитного, существует
также радиальное электрическое поле напряжённости $E = 2e/lr^2$, где
$e$ — заряд на внутренней обкладке конденсатора, численно равный, но
противоположный по знаку заряду на его внешней обкладке, $l$ — длина
конденсатора, а $r$ — векторное расстояние точки поля от оси конденсатора
[см.\ уравнение (4.5)].

Мы предполагаем, что длина конденсатора настолько больше его диаметра,
что краевыми эффектами можно пренебречь, т.\,е. можно пользоваться
формулой (4.5), строго справедливой лишь для бесконечно длинного
конденсатора. Далее, для простоты предполагаем, что в полости конденсатора
находится газ, отличием диэлектрической постоянной которого от единицы
можно пренебречь.

В пространстве между обкладками конденсатора вектор Пойнтинга отличен
от нуля и равен:
$$
    \mathbf{S} = \frac{c}{4\pi}[\mathbf{E}\mathbf{H}]
               = \frac{ce}{2\pi r^2 l}[\mathbf{r}\mathbf{H}].
$$

% стр. 503 (правая колонка)

Линии вектора Пойнтинга, т.\,е. линии потока энергии, представляют
собой концентрические окружности, плоскости которых перпендикулярны
оси конденсатора.

Таким образом, мы приходим к представлению о беспрерывной циркуляции
энергии по замкнутым путям в \textit{статическом} электромагнитном поле.
Представление это не приводит к каким-либо следствиям, могущим быть
непосредственно проверенными на опыте, а потому лишено непосредственного
физического смысла.

Примем, однако, во внимание, что, согласно (103.8), плотность
электромагнитного количества движения $\mathbf{g}$ пропорциональна
вектору Пойнтинга $\mathbf{S}$\,\footnote{Заметим, что соотношение
$\mathbf{g} = \dfrac{1}{c^2}\mathbf{S}$ вытекает из основных положений
теории относительности и должно выполняться для любых полей в вакууме,
а также и для всех вообще форм энергии и импульса при строгой
микроскопической их трактовке.}. Утверждение, что в рассматриваемом
статическом поле локализовано определённое количество движения с
объёмной плотностью
\begin{equation}
    \mathbf{g} = \frac{1}{c^2}\mathbf{S} = \frac{e}{2\pi cr^2 l}[\mathbf{r}\mathbf{H}],
    \tag{104.1}
\end{equation}
является содержательным высказыванием и приводит к следствиям, доступным
(по крайней мере принципиально) опытной проверке.

Правда, общее количество движения всего \textit{статического} поля
в целом по необходимости равно нулю. Однако, зная пространственное
распределение электромагнитного количества движения, мы можем определить
\textit{момент} этого количества движения $\mathbf{K}$ относительно
центра инерции конденсатора по формуле
\begin{equation}
    \mathbf{K} = \int[\mathbf{R}\mathbf{g}]\,dV,
    \tag{104.2}
\end{equation}
которая совершенно аналогична формуле
$$
    \mathbf{K} = \sum_i[\mathbf{R}_i\mathbf{g}_i] = \sum_i m_i[\mathbf{R}_i\mathbf{v}_i],
$$
определяющей момент механического количества движения системы
материальных точек масс $m_i$ и скорости $\mathbf{v}_i$. Если мы разрядим
конденсатор, то исчезнет как электрическое поле $\mathbf{E}$, так и
электромагнитный момент количества движения $\mathbf{K}$. На основании
закона сохранения момента количества движения мы должны заключить, что
система конденсатор $+$ магнит, возбуждающий поле $\mathbf{H}$, приобретает
в процессе разряда равный $\mathbf{K}$ момент \textit{механического}
количества движения. Если, например, магнит закреплён неподвижно,
а конденсатор может свободно вращаться около своей оси, то в

% стр. 504 (левая колонка)

процессе разряда он должен приобрести угловую скорость вращения, равную
$\omega = K/I$, где $I$ — момент инерции конденсатора относительно его оси.

Этот доступный опытной проверке вывод, вытекающий из предположения
о локализации в электромагнитном поле количества движения с плотностью
$\mathbf{g}$, может быть подтверждён непосредственным расчётом, к которому
мы теперь и перейдём.

\textbf{2.} Применение закона сохранения момента количества движения
предполагает, что рассматриваемая система не подвергается во время
разряда действию внешних сил (либо что сумма этих сил равна нулю). Это
условие будет удовлетворено, если мы предположим, например, что разряд
конденсатора производится путём приближения к нему радиоактивного
вещества, вызывающего ионизацию газа между обкладками конденсатора.
Если плотность разрядного тока в газе равна $\mathbf{j}$, то каждый
элемент объёма $dV$ будет во время разряда испытывать силу
$\dfrac{1}{c}[\mathbf{j}\mathbf{H}]\,dV$, момент же всех сил, испытываемых
газом, будет равен:
\begin{equation}
    \mathbf{N} = \frac{1}{c}\int[\mathbf{R}[\mathbf{j}\mathbf{H}]]\,dV,
    \tag{104.3}
\end{equation}
где интегрирование должно быть распространено на всё пространство
между обкладками конденсатора. Благодаря трению между газом и обкладками
этот момент будет в конечном счёте передан всему конденсатору в целом.

Введём цилиндрическую систему координат $z$, $r$, $\alpha$, ось $r$
которой совпадает с осью конденсатора. Если бы магнитного поля не было,
то вектор плотности тока был бы радиальным ($\mathbf{j} = |j_r|$).
В магнитном поле $\mathbf{H}$ линии тока благодаря эффекту Холла
приобретут спиралевидную форму, и слагающая $j_\alpha$ будет отличной
от нуля, слагающая же $j_z$ ввиду аксиальной симметрии поля останется
равной нулю: $j_z = 0$.

Далее, ввиду той же симметрии момент $\mathbf{N}$ приложенных к газу
сил будет направлен по оси $z$
$$
    (N_x = N_y = 0).
$$

Раскрывая тройное векторное произведение в формуле (104.3), получим:
$$
    N_z = \frac{1}{c}\int\{j(RH) - H(Rj)\}\,dV.
$$

Приняв во внимание, что поле $\mathbf{H}$ однородно, что $j_z = 0$
и что поэтому $\mathbf{R}\mathbf{j} = r\mathbf{j} = rj_r$, получаем:
$$
    N_z = -\frac{H_z}{c}\int rj_r\,dV.
$$

% стр. 505 (правая колонка)

Ток, протекающий в момент времени $t$ через коаксиальную с конденсатором
цилиндрическую поверхность радиуса $r$, равен:
$$
    J(t) = 2\pi r l\cdot j_r,
$$
причём ток считается положительным, если он направлен от оси конденсатора
наружу. Сила тока $J(t)$ не зависит от радиуса $r$ цилиндрической
поверхности (мы пренебрегаем возможностью накопления объёмных зарядов
в полости конденсатора\,\footnote{Если допустить образование объёмных
зарядов, то соответственно нужно видоизменить выражение для напряжённости
электрического поля в конденсаторе. В результате вновь придём к формуле
(104.5).}). Поэтому
$$
    N_z = -\frac{H_z J(t)\cdot V}{2\pi cl},
$$
где $V$ — объём полости конденсатора. Наконец, сила тока $J(t)$,
протекающего через конденсатор от его внутренней обкладки к внешней,
равна быстроте убыли заряда $e$ на внутренней обкладке конденсатора:
$$
    J(t) = -\frac{de}{dt}.
$$

Так как векторы $\mathbf{N}$ и $\mathbf{H}$ направлены по оси $Z$, то
окончательно
$$
    N_z = \frac{HV}{2\pi cl}\frac{de}{dt}.
$$

Механический момент количества движения конденсатора $\mathbf{K}_M$
будет под воздействием момента сил $\mathbf{N}$ изменяться по закону
\begin{equation}
    \frac{d}{dt}\mathbf{K}_M = \mathbf{N} = \frac{HV}{2\pi cl}\frac{de}{dt}.
    \tag{104.4}
\end{equation}

С другой стороны, общий момент количества движения поля внутри
конденсатора, согласно (104.2) и (104.1), равен:
$$
    \mathbf{K} = \int[\mathbf{R}\mathbf{g}]\,dV
              = \frac{e}{2\pi cl}\int\frac{[\mathbf{R}[\mathbf{r}\mathbf{H}]]}{r^2}\,dV.
$$

Тройное произведение под знаком интеграла равно:
$$
    [\mathbf{R}[\mathbf{r}\mathbf{H}]] = \mathbf{r}(\mathbf{R}\mathbf{H})
    - \mathbf{H}(\mathbf{R}\mathbf{r}),
$$
причём
$$
    \mathbf{R}\mathbf{r} = r^2.
$$
Поэтому
$$
    K_z = -\frac{eH_z V}{2\pi cl}.
$$

% стр. 506 (левая колонка)

Так как из соображений симметрии явствует, что слагающие по осям
$x$ и $y$ вектора $\mathbf{K}$ равны нулю, то, сравнивая последнее
уравнение с (104.4), получаем:
$$
    \frac{d}{dt}\mathbf{K}_M = -\frac{d}{dt}\mathbf{K},
$$
или окончательно
\begin{equation}
    \frac{d}{dt}(\mathbf{K} + \mathbf{K}_M) = 0,
    \tag{104.5}
\end{equation}

Таким образом, мы подтвердили непосредственным подсчётом, что сумма
электромагнитного и механического моментов количества движения остаётся
постоянной во времени, так что при разряде конденсатора в магнитном поле
он должен приобрести механический момент количества движения $\mathbf{K}_M$.

\textbf{3.} Рассмотренный пример показывает, что понятие электромагнитного
момента количества движения оказывается плодотворным даже в случае
статических полей. Ещё большую роль играет это понятие при изучении
переменных полей, в особенности поля излучения.

Правда, в области макроскопических явлений экспериментальное измерение
электромагнитного количества движения и его момента представляет очень
трудную задачу ввиду ничтожной величины связанных с ними эффектов\,\footnote{Ср.\
сказанное в §\,103 о давлении света. Экспериментально удалось измерить
момент количества движения, передаваемый светом кварцевой пластинке при
прохождении его через эту пластинку; передача момента связана с вращением
плоскости поляризации света в пластинке. Этот эффект, обнаруженный А.\,Бэтом
(\textit{Phys. Rev.} 50, 115, 1936), был предсказан теоретически русским
учёным А.\,Садовским ещё в 1899 г.}. Однако в области атомных явлений
обмен моментом количества движения между светом и веществом имеет весьма
существенное значение. Так, например, излучение света возбуждённым
атомом, вообще говоря, связано с изменением момента количества движения
электронов в атоме, причём это изменение по порядку величины сравнимо
с абсолютной величиной момента атома.

% ============================================================
% §105. ТЕНЗОР НАТЯЖЕНИЙ И ПОНДЕРОМОТОРНЫЕ СИЛЫ
% ЭЛЕКТРОМАГНИТНОГО ПОЛЯ (стр. 506–511)
% ============================================================

\section*{\S\,105. Тензор натяжений и пондеромоторные силы
электромагнитного поля.}

% стр. 506 (левая колонка, продолжение)

\textbf{1.} В §§\,34 и 84 мы нашли выражение тензора натяжений
$\mathbf{T}$ в стационарных электрическом и магнитном полях. Покажем
теперь, что если допустить применимость этих выражений и в случае
произвольного переменного электромагнитного поля, то вытекающие из
этого допущения выводы будут находиться в полном соответствии с
результатами настоящей главы.

% стр. 507 (правая колонка)

В §§\,34 и 84 мы убедились, что тензор напряжений электромагнитного
поля $\mathbf{T}$ может быть разложен на сумму «максвеллова» тензора
$\mathbf{T}'$ и тензора стрикционных натяжений $\mathbf{T}''$. В
дальнейшем мы для простоты вовсе не будем рассматривать стрикционных
натяжений $\mathbf{T}''$, приводящих лишь к перераспределению
пондеромоторных сил по объёму находящихся в поле тел, но не изменяющих
результирующей этих сил, и будем отождествлять полный тензор натяжений
$\mathbf{T}$ с максвелловым тензором $\mathbf{T}'$. При этих условиях
тензор натяжений $\mathbf{T}$ выразится суммой выражений (34.2) и (84.4),
которую можно представить в следующей форме:
\begin{equation}
    T_{ik} = \frac{1}{8\pi}\{E_i D_k + E_k D_i + H_i B_k + H_k B_i
           - \delta_{ik}(\mathbf{ED} + \mathbf{HB})\}.
    \tag{105.1}
\end{equation}

Индексы $i$ и $k$ в этом выражении могут пробегать значения $x$, $y$,
$z$, а величина $\delta_{ik}$ определяется уравнениями
\begin{equation}
    \delta_{ii} = \delta_{xx} = \delta_{yy} = \delta_{zz} = 1,
    \quad
    \delta_{ik} = 0 \text{ при } i \neq k.
    \tag{105.2}
\end{equation}

Впрочем, удобнее ввести обозначения $x = x_1$, $y = x_2$, $z = x_3$
и придавать индексам $i$ и $k$ значения 1, 2, 3. В этих обозначениях
формула (33.7), выражающая плотность пондеромоторных сил $\mathbf{f}$,
принимает вид
$$
    f_i = \sum_{k=1}^3 \frac{\partial T_{ik}}{\partial x_k}.
$$

Воспользовавшись значениями (105.1) слагающих $T_{ik}$, которые до
дифференцирования удобно представить в виде
$$
    T_{ik} = \frac{1}{4\pi}(E_i D_k + H_i B_k)
           - \frac{1}{8\pi}\delta_{ik}(\mathbf{ED} + \mathbf{HB}),
$$
получаем:
$$
    \sum_k\frac{\partial T_{ik}}{\partial x_k}
    = \frac{1}{4\pi}\sum_k\!\left(
        E_i\frac{\partial D_k}{\partial x_k}
      + \frac{\partial E_i}{\partial x_k}D_k
      + H_i\frac{\partial B_k}{\partial x_k}
      + \frac{\partial H_i}{\partial x_k}B_k
    \right)
    - \frac{1}{8\pi}\frac{\partial}{\partial x_i}(\mathbf{ED} + \mathbf{HB}).
$$

Последний член получается на основании (105.2):
$$
    -\frac{1}{8\pi}\sum_k\delta_{ik}\frac{\partial}{\partial x_k}(\mathbf{ED} + \mathbf{HB})
    = -\frac{1}{8\pi}\delta_{ii}\frac{\partial}{\partial x_i}(\mathbf{ED} + \mathbf{HB})
    = -\frac{1}{8\pi}\frac{\partial}{\partial x_i}(\mathbf{ED} + \mathbf{HB}).
$$

% стр. 508 (левая колонка)

Примем во внимание, что, согласно максвелловым уравнениям (III) и (IV),
$$
    \sum_k\frac{\partial D_k}{\partial x_k}
    = \frac{\partial D_x}{\partial x} + \frac{\partial D_y}{\partial y}
    + \frac{\partial D_z}{\partial z} = \operatorname{div}\mathbf{D} = 4\pi\rho,
$$
$$
    \sum_k\frac{\partial B_k}{\partial x_k} = \operatorname{div}\mathbf{B} = 0.
$$

Далее,
$$
    \frac{\partial}{\partial x_i}(\mathbf{ED})
    = \sum_k\frac{\partial}{\partial x_i}(E_k D_k)
    = \sum_k\!\left(\frac{\partial E_k}{\partial x_i}D_k
    + E_k\frac{\partial D_k}{\partial x_i}\right)
$$
и, стало быть,
\begin{align*}
    \frac{1}{4\pi}\sum_k\frac{\partial E_i}{\partial x_k}D_k
    - \frac{1}{8\pi}\frac{\partial}{\partial x_i}(\mathbf{ED})
    &= \frac{1}{4\pi}\sum_k\!\left(\frac{\partial E_i}{\partial x_k}
     - \frac{\partial E_k}{\partial x_i}\right)D_k
     + \frac{1}{8\pi}\sum_k\!\left(\frac{\partial E_k}{\partial x_i}D_k
     - E_k\frac{\partial D_k}{\partial x_i}\right).
\end{align*}

Аналогичным образом можно преобразовать и члены
$$
    \frac{1}{4\pi}\sum_k\frac{\partial H_i}{\partial x_k}B_k
    - \frac{1}{8\pi}\frac{\partial}{\partial x_i}(\mathbf{HB}).
$$

В рассматриваемом нами случае неподвижной изотропной среды связь
между $\mathbf{E}$ и $\mathbf{D}$ и между $\mathbf{H}$ и $\mathbf{B}$
определяется уравнениями Максвелла (V), и поэтому
$$
    \sum_k\!\left(\frac{\partial E_k}{\partial x_i}D_k
    - E_k\frac{\partial D_k}{\partial x_i}\right)
    = -\sum_k E_k\frac{\partial\varepsilon}{\partial x_i}
    = -E^2\frac{\partial\varepsilon}{\partial x_i},
$$
$$
    \sum_k\!\left(\frac{\partial H_k}{\partial x_i}B_k
    - H_k\frac{\partial B_k}{\partial x_i}\right)
    = -H^2\frac{\partial\mu}{\partial x_i}.
$$

Поэтому $\sum_k\dfrac{\partial T_{ik}}{\partial x_k}$ может быть
представлена следующим образом:
\begin{equation}
    \sum_k\frac{\partial T_{ik}}{\partial x_k}
    = \rho E_i + \frac{1}{4\pi}\sum_k\!\left(
        \frac{\partial E_i}{\partial x_k} - \frac{\partial E_k}{\partial x_i}
    \right)D_k
    + \frac{1}{4\pi}\sum_k\!\left(
        \frac{\partial H_i}{\partial x_k} - \frac{\partial H_k}{\partial x_i}
    \right)B_k
    - \frac{1}{8\pi}\!\left(
        E^2\frac{\partial\varepsilon}{\partial x_i}
      + H^2\frac{\partial\mu}{\partial x_i}
    \right).
    \tag{105.3}
\end{equation}

Рассмотрим слагающую выражения (105.3) по оси $x$. Второй член справа
при $i = 1$ равен:
$$
    \frac{1}{4\pi}\sum_{k=1}^3\!\left(
        \frac{\partial E_1}{\partial x_k} - \frac{\partial E_k}{\partial x_1}
    \right)D_k
    = \frac{1}{4\pi}\!\left(
        \frac{\partial E_1}{\partial x_2} - \frac{\partial E_2}{\partial x_1}
    \right)D_2
    + \frac{1}{4\pi}\!\left(
        \frac{\partial E_1}{\partial x_3} - \frac{\partial E_3}{\partial x_1}
    \right)D_3.
$$

% стр. 509 (правая колонка)

В обычных обозначениях
$$
    \frac{\partial E_1}{\partial x_2} - \frac{\partial E_2}{\partial x_1}
    = \frac{\partial E_x}{\partial y} - \frac{\partial E_y}{\partial x}
    = -\operatorname{rot}_z\mathbf{E}
$$
и
$$
    \frac{\partial E_1}{\partial x_3} - \frac{\partial E_3}{\partial x_1}
    = \operatorname{rot}_y\mathbf{E}.
$$

Следовательно,
$$
    \sum_{k=1}^3\!\left(
        \frac{\partial E_i}{\partial x_k} - \frac{\partial E_k}{\partial x_i}
    \right)D_k
    = -D_y\operatorname{rot}_z\mathbf{E} + D_z\operatorname{rot}_y\mathbf{E}
    = [\operatorname{rot}\mathbf{E}\cdot\mathbf{D}]_x
    = [\operatorname{rot}\mathbf{E}\cdot\mathbf{D}]_i.
$$

Аналогично этому может быть преобразован и третий член справа в (105.3),
так что сумма второго и третьего членов равна:
$$
    \frac{1}{4\pi}\sum_k\!\left(
        \frac{\partial E_i}{\partial x_k} - \frac{\partial E_k}{\partial x_i}
    \right)D_k
    + \frac{1}{4\pi}\sum_k\!\left(
        \frac{\partial H_i}{\partial x_k} - \frac{\partial H_k}{\partial x_i}
    \right)B_k
    = \frac{1}{4\pi}[\operatorname{rot}\mathbf{E}\cdot\mathbf{D}]_i
    + \frac{1}{4\pi}[\operatorname{rot}\mathbf{H}\cdot\mathbf{B}]_i.
$$

На основании максвелловых уравнений (I) и (II)
$$
    [\operatorname{rot}\mathbf{E}\cdot\mathbf{D}]
    + [\operatorname{rot}\mathbf{H}\cdot\mathbf{B}]
    = -\frac{1}{c}\!\left[\frac{\partial\mathbf{D}}{\partial t}\cdot\mathbf{D}\right]
    + \frac{1}{c}\!\left[\frac{\partial\mathbf{B}}{\partial t}\cdot\mathbf{B}\right]
    + \frac{4\pi}{c}[\mathbf{j}\mathbf{B}].
$$

Далее,
$$
    -\frac{1}{c}\!\left[\frac{\partial\mathbf{D}}{\partial t}\cdot\mathbf{D}\right]
    + \frac{1}{c}\!\left[\frac{\partial\mathbf{B}}{\partial t}\cdot\mathbf{B}\right]
    = \frac{1}{c}\!\left[\mathbf{D}\frac{\partial\mathbf{B}}{\partial t}\right]
    + \frac{1}{c}\!\left[\frac{\partial\mathbf{D}}{\partial t}\mathbf{B}\right]
    = \frac{1}{c}\frac{\partial}{\partial t}[\mathbf{D}\mathbf{B}].
$$

Таким образом, окончательно получаем:
\begin{equation}
    \sum_k\frac{\partial T_{ik}}{\partial x_k}
    = \rho E_i + \frac{1}{c}[\mathbf{j}\mathbf{B}]_i
    - \frac{1}{8\pi}\!\left(
        E^2\frac{\partial\varepsilon}{\partial x_i}
      + H^2\frac{\partial\mu}{\partial x_i}
    \right)
    + \frac{\partial}{\partial t}\!\left(\frac{1}{4\pi c}[\mathbf{D}\mathbf{B}]\right)_i.
    \tag{105.4}
\end{equation}

В рассматриваемом нами случае неподвижной изотропной среды последний
член в уравнении (105.4) может быть представлен в следующей форме:
$$
    \frac{\partial}{\partial t}\!\left(\frac{1}{4\pi c}[\mathbf{D}\mathbf{B}]\right)_i
    = \frac{\partial g_i}{\partial t}
    + \frac{(\varepsilon\mu - 1)}{4\pi c}\frac{\partial}{\partial t}[\mathbf{E}\mathbf{H}],
$$
где вектор
$$
    \mathbf{g} = \frac{1}{4\pi c}[\mathbf{E}\mathbf{H}]
$$

% стр. 510 (левая колонка)

означает, согласно уравнению (103.8), плотность количества движения
электромагнитного поля.

Если ввести ещё обозначения
\begin{equation}
    f_i^{(1)} = \rho E_i + \frac{1}{c}[\mathbf{j}\mathbf{B}]_i,
    \tag{105.5}
\end{equation}
\begin{equation}
    f_i^{(2)} = -\frac{1}{8\pi}\!\left(
        E^2\operatorname{grad}\varepsilon + H^2\operatorname{grad}\mu
    \right)_i
    + \frac{(\varepsilon\mu - 1)}{4\pi c}\frac{\partial}{\partial t}[\mathbf{E}\mathbf{H}]_i,
    \tag{105.6}
\end{equation}
то уравнение (105.4) примет вид
\begin{equation}
    \sum_k\frac{\partial T_{ik}}{\partial x_k}
    = f_i^{(1)} + f_i^{(2)} + \frac{\partial g_i}{\partial t}.
    \tag{105.7}
\end{equation}

Первый член справа $f^{(1)}$ равен слагающей по оси $i$ плотности
пондеромоторных сил, действующих на свободные заряды $\rho$ и на
электрические токи $\mathbf{j}$. Второй член справа в (105.4)
$\mathbf{f}^{(2)}$ представляет собой плотность пондеромоторных сил,
действующих на среду (т.\,е. на диэлектрики и магнетики). Действительно,
в вакууме $\varepsilon = \mu = 1$ и поэтому, согласно (105.6), в вакууме
$\mathbf{f}^{(2)} = 0$.

Кроме того, в постоянном поле последний член в уравнении (105.6)
равен нулю и определённый этим уравнением вектор действительно совпадает
с суммой вторых членов формул (32.10) и (83.5) для плотности сил,
действующих на диэлектрики и магнетики.

\textbf{2.} Если обозначить через $\mathbf{f}$ полную плотность
пондеромоторных сил:
\begin{equation}
    \mathbf{f} = \mathbf{f}^{(1)} + \mathbf{f}^{(2)},
    \tag{105.8}
\end{equation}
где $\mathbf{f}^{(1)}$ и $\mathbf{f}^{(2)}$ определяются уравнениями
(105.5) и (105.6), то уравнение (105.4) примет вид
\begin{equation}
    f_i + \frac{\partial g_i}{\partial t}
    = \sum_k\frac{\partial T_{ik}}{\partial x_k}.
    \tag{105.9}
\end{equation}

В стационарных полях $\dfrac{\partial g_i}{\partial t} = 0$ и это
уравнение совпадает с (33.7); в переменных полях оно отличается от (33.7)
членом $\dfrac{\partial g_i}{\partial t}$, учитывающим изменение во
времени электромагнитного количества движения $\mathbf{g}$. Этот
дополнительный член обеспечивает сохранение суммы количества движения
материальных тел и электромагнитного поля.

Действительно, проинтегрируем уравнение (105.9) по всему пространству,
предполагая, что слагающие электромагнитного тензора натяжений достаточно
быстро убывают по мере удаления в бесконечность. Интеграл правой части
уравнения (105.9) при этих усло-

% стр. 511 (правая колонка)

виях обратится в нуль, ибо, например, при $i = 1$
$$
    \sum_k\frac{\partial T_{ik}}{\partial x_k}
    = \frac{\partial T_{xx}}{\partial x}
    + \frac{\partial T_{xy}}{\partial y}
    + \frac{\partial T_{xz}}{\partial z},
$$
и, стало быть, согласно уравнению, непосредственно предшествующему
уравнению (33.7) на стр. 159:
$$
    \int_V\sum_k\frac{\partial T_{ik}}{\partial x_k}\,dV
    = \oint_S T_{xn}\,dS,
$$
где поверхностный интеграл должен быть взят по поверхности $S$ объёма
$V$ и при удалении этой поверхности в бесконечность обратиться, по
предположению, в нуль. Таким образом, получаем:
$$
    \int f_i\,dV + \frac{\partial}{\partial t}\int g_i\,dV = 0,
$$
или, переходя от слагающих к векторам и заменяя частную производную
по времени полной (ввиду того, что под знаком $\dfrac{\partial}{\partial t}$
стоит интеграл по всему пространству, зависящий только от $t$):
$$
    \frac{d}{dt}\int\mathbf{g}\,dV + \int\mathbf{f}\,dV = 0.
$$

С другой стороны, изменение механического количества движения
$\mathbf{G}_M$ всех находящихся в пространстве тел определяется
равнодействующей всех пондеромоторных сил:
$$
    \frac{d\mathbf{G}_M}{dt} = \int\mathbf{f}\,dV.
$$

Поэтому последнее равенство равносильно равенству
\begin{equation}
    \frac{d}{dt}\!\left(\mathbf{G} + \int\mathbf{g}\,dV\right) = 0,
    \tag{105.10}
\end{equation}
которое и выражает собой \textit{закон сохранения полного количества
движения} (механического плюс электромагнитного).

Важное уравнение (105.9) устанавливает в самой общей форме связь
между плотностью пондеромоторных сил $\mathbf{f}$, количеством движения
поля $\mathbf{g}$ и натяжениями $T_{ik}$.

% ============================================================
% §106. ПРИМЕР НЕКВАЗИСТАЦИОНАРНЫХ ТОКОВ: ВОЛНЫ ВДОЛЬ КАБЕЛЯ
% (стр. 511 — начало)
% ============================================================

% \section*{\S\,106. Пример неквазистационарных токов: волны вдоль кабеля.}

% \textbf{1.} В главе VI мы подробно рассмотрели свойства квазистационарных
% переменных токов. Теперь же мы рассмотрим на частном примере свойства
% \textit{быстропеременных} токов, к которым теория, изложенная

















\end{document}
